\BourbakiExerciseGroup

\begin{enumerate}
\def\labelenumi{\arabic{enumi})}
\item
  设 \(\mathcal{C}_1, \mathcal{C}_2\) 是同一集合 X 上的两个拓扑。证明：\(\mathcal{C}_1\) 比 \(\mathcal{C}_2\) 更细，当且仅当 X 上每个在拓扑 \(\mathcal{C}_1\) 中收敛的滤子都在拓扑 \(\mathcal{C}_2\) 中收敛到同一点。
\item
  设 \(\mathfrak{U}\) 是 \(\mathbf{N}\) 上比 Fréchet 滤子更细的一个超滤子，并设 \(\mathbf{X} = \mathbf{N} \cup \{\omega\}\) 为与 \(\mathfrak{U}\) 相伴随的空间（§ 6，第 5 目）。证明：具有无穷多个不同项的序列 \((x_n)\) 在 \(\mathbf{X}\) 中决不收敛。
\item
  设 X, \(X'\) 为两个拓扑空间，\(f: X \to X'\) 为在点 \(x_{0} \in X\) 处连续的映射。证明：若 B 是 X 上以 \(x_{0}\) 为聚点的任一滤子基，则 \(f(x_{0})\) 是滤子基 \(f(\mathfrak{B})\) 在 \(X'\) 上的聚点。
\item
  \begin{enumerate}
  \def\labelenumii{\alph{enumii})}
  \tightlist
  \item
    设 X, Y 为两个拓扑空间，F 为 X 上以 a 为聚点的滤子，G 为 Y 上以 b 为聚点的滤子。证明 \((a, b)\) 是积滤子 \(F \times G\) 的聚点。
  \end{enumerate}
\end{enumerate}

\begin{enumerate}
\def\labelenumi{\alph{enumi})}
\setcounter{enumi}{1}
\tightlist
\item
  在积空间 \(Q^{2}\) 中，给出一个序列 \((x_{n}, y_{n})\) 的例子：它没有聚点，尽管序列 \((x_{n})\) 与 \((y_{n})\) 各自在 Q 中都有聚点。
\end{enumerate}

\begin{enumerate}
\def\labelenumi{\arabic{enumi})}
\setcounter{enumi}{4}
\item
  设 X, Y 为两个拓扑空间，G 是 \(X \times Y\) 中映射 \(f: X \to Y\) 的图像。证明：对每个 \(x \in X\)，f 在 x 处的聚点所成的集就是 \(\overline{\mathbf{G}}(x)\)，即 \(\overline{\mathbf{G}}\)（G 在 \(X \times Y\) 中的闭包）在 x 处的截口。
\item
  设 \((\mathbf{X}_t)_{t\in I}\) 是由离散空间组成的一个不可数族，其中每个空间至少有两个点。在积集合上考虑拓扑
\end{enumerate}

\[
\mathbf {X} = \prod_ {\iota \in \mathbf {I}} \mathbf {X} _ {\iota},
\]

它由所有的积 \(\prod_{i\in I}M_i\) 生成，其中 \(M_i\subset X_i\) 对一切 \(i\in I\) 成立，且除去可数个指标外均有 \(M_i = X_i\)（参见§ 4，习题 9）。证明：开集的每个可数交在 \(X\) 中都是开集，并由此推出，\(x\) 的任何点都不具有可数的基本邻域系。

\begin{enumerate}
\def\labelenumi{\arabic{enumi})}
\setcounter{enumi}{6}
\tightlist
\item
  拓扑空间 X 的子集 A 称为本原的，如果它是 X 上某个超滤子的极限点集。设 \(Y_{A}\) 表示 X 上以 A 为极限点集的超滤子所成的集。
\end{enumerate}

\begin{enumerate}
\def\labelenumi{\alph{enumi})}
\item
  若 A 是 X 的本原子集，证明 X 的每个与 A 相交的开子集都属于所有超滤子 \(U \in Y_{A}\)。
\item
  设 A, B 是 X 的两个不同的本原子集。证明存在超滤子 \(U \in Y_{A}\)、超滤子 \(B \in Y_{B}\) 以及 X 的子集 M，使得 \(M \in U\) 且 \(CM \in B\)。
\item
  设 \(f: \mathbf{X} \to \mathbf{Y}\) 是连续映射。证明在 \(f\) 下，\(\mathbf{X}\) 的任何本原子集的像都含于 \(\mathbf{Y}\) 的某个本原子集。
\item
  若 \(x \in \mathbf{X}\) 且 \(\{x\}\) 在 \(\mathbf{X}\) 中闭，证明 \(\{x\}\) 是 \(\mathbf{X}\) 的本原子集。
\end{enumerate}

\BourbakiExerciseGroup

\begin{enumerate}
\def\labelenumi{\arabic{enumi})}
\tightlist
\item
  \begin{enumerate}
  \def\labelenumii{\alph{enumii})}
  \tightlist
  \item
    设 X 是拓扑空间。证明以下三个命题等价：
  \end{enumerate}
\end{enumerate}

\begin{enumerate}
\def\labelenumi{(\Alph{enumi})}
\setcounter{enumi}{16}
\tightlist
\item
  若 \(x, y\) 是 \(X\) 的任意两个不同的点，则存在 \(x\) 的一个邻域，它不含 \(y\)。
\end{enumerate}

(Q') 由单独一点组成的 X 的每个子集在 X 中闭。(Q'\,') 对每个 \(x \in X\)，x 的诸邻域之交仅由点 x 组成。

满足这些条件的空间 X 称为可达的。

\begin{enumerate}
\def\labelenumi{\alph{enumi})}
\setcounter{enumi}{1}
\tightlist
\item
  带右拓扑的有序集 X（关于这个拓扑 X 是 Kolmogoroff 空间{[}§ 1，习题 2{]}）是可达的，当且仅当 X 的任何两个不同元素都不可比较，且此时 X 上的右拓扑是离散拓扑。
\end{enumerate}

\begin{enumerate}
\def\labelenumi{\arabic{enumi})}
\setcounter{enumi}{1}
\tightlist
\item
  \begin{enumerate}
  \def\labelenumii{\alph{enumii})}
  \tightlist
  \item
    若集合 X 上的一个拓扑具有有限子基，且 X 关于这个拓扑是可达的，则 X 有限且该拓扑是离散拓扑。
  \end{enumerate}
\end{enumerate}

\begin{enumerate}
\def\labelenumi{\alph{enumi})}
\setcounter{enumi}{1}
\tightlist
\item
  若在可达空间 X 中开集的每个交都是开集，则 X 是离散的（参见§ 7，习题 6）。
\end{enumerate}

\begin{enumerate}
\def\labelenumi{\arabic{enumi})}
\setcounter{enumi}{2}
\tightlist
\item
  \begin{enumerate}
  \def\labelenumii{\alph{enumii})}
  \tightlist
  \item
    设 X 是可达空间，A 是 X 的子集，x 是 \(\overline{A}\) 中不属于 A 的一点。证明 x 的每个邻域都含有 A 的无穷多个点。
  \end{enumerate}
\end{enumerate}

\begin{enumerate}
\def\labelenumi{\alph{enumi})}
\setcounter{enumi}{1}
\item
  设 X 是可达空间，A 是 X 的子集。证明由这样的点 \(x \in \overline{A}\) 组成的集在 X 中闭：x 的每个邻域都含有 A 的异于 x 的点。
\item
  设 X 是可达空间。证明 X 的任何子集 A 的一切邻域之交等于 A。
\end{enumerate}

\begin{enumerate}
\def\labelenumi{\arabic{enumi})}
\setcounter{enumi}{3}
\tightlist
\item
  证明 Kolmogoroff 空间（相应地，可达空间）的每个子空间是 Kolmogoroff 空间（相应地，可达空间）；并且比 Kolmogoroff 空间（相应地，可达空间）的某个拓扑更细的拓扑，仍是 Kolmogoroff 空间（相应地，可达空间）的拓扑。
\end{enumerate}

¶ 5) a) 设 X 是无限集。证明 X 上所有拓扑 \(\mathcal{G}_{\mathfrak{c}}\)（§ 2，习题 8），只要 \(1 < c \leqslant\) card (X)，就都是可达空间的拓扑，但都不是豪斯多夫的。

\begin{enumerate}
\def\labelenumi{\alph{enumi})}
\setcounter{enumi}{1}
\tightlist
\item
  证明 X 上所有豪斯多夫拓扑之交是非豪斯多夫拓扑 \(\mathfrak{G}_{\mathrm{card}(\mathbf{N})}\)，它也是 X 上最粗的可达空间拓扑。（利用§ 2 的习题 8，并注意 X 上存在这样的豪斯多夫拓扑：关于它，X 有非闭的可数无限子集。）
\end{enumerate}

\begin{enumerate}
\def\labelenumi{\arabic{enumi})}
\setcounter{enumi}{5}
\tightlist
\item
  \begin{enumerate}
  \def\labelenumii{\alph{enumii})}
  \tightlist
  \item
    设 X 是豪斯多夫空间，D 是 X 的稠密子集。证明 \(\text{card}(X) \leqslant 2^{2^{\text{card}(D)}}\)（注意 X 的每一点都是 D 上某个滤子基的极限）。证明 card (X) 的这个上界不能改进（§ 4，习题 5 b)），并由此推出：在无限集 A 上，超滤子所成的集与滤子所成的集都和 \(\mathfrak{P}(\mathfrak{P}(A))\) 等势。
  \end{enumerate}
\end{enumerate}

\begin{enumerate}
\def\labelenumi{\alph{enumi})}
\setcounter{enumi}{1}
\tightlist
\item
  设 K 是由数 o 和 1 组成的离散空间，A 是无限集。由 a) 推出：若 card (A) ≥ 2 \(^{2\text{card(N)}}\)，则积空间 K \(^{\text{A}}\) 满足§ 1 习题 7 的条件 (D \(_{\text{IV}}\) )，但既不满足 (D \(_{\text{II}}\) ) 也不满足 (D \(_{\text{III}}\) )。（§ 4，习题 4。）
\end{enumerate}

* 7) 设 X 是 R 中的区间 \([-1, 1]\)。考虑 X 上如下的等价关系 S：若 \(x \neq \pm 1\)，则 \(x\) 的等价类由 \(x\) 和 \(-x\) 组成；\(1\)（相应地，\(-1\)）的等价类仅由 \(1\)（相应地，\(-1\)）组成。证明 S 是开的，且商空间 X/S 可达但非豪斯多夫。再证明 X/S 的每一点在 X/S 中都有一个作为正则子空间的邻域。*

\begin{enumerate}
\def\labelenumi{\arabic{enumi})}
\setcounter{enumi}{7}
\tightlist
\item
  \begin{enumerate}
  \def\labelenumii{\alph{enumii})}
  \tightlist
  \item
    设 X 是子空间族 \((\mathbf{X}_{t})\) 的和所成的豪斯多夫（相应地，正则）空间，X/R 是把诸 \(X_{t}\) 沿闭子集 \(A_{tx}\) 借助同胚 \(h_{x_{t}}\)（§ 3，第 4 目）粘合起来而得的商空间；再设对每个指标 t，使得 \(A_{tx} \neq \emptyset\) 的指标 x 只有有限个。在这些条件下，证明 X/R 是豪斯多夫（相应地，正则）的。
  \end{enumerate}
\end{enumerate}

* b) 证明习题 7 的空间 X/S 可以通过沿两个开集把两个正则子空间粘合在一起而得到。*

\begin{enumerate}
\def\labelenumi{\arabic{enumi})}
\setcounter{enumi}{8}
\tightlist
\item
  设 \(\Gamma\) 是豪斯多夫（相应地，正则）空间 X 的一个有限同胚群，R 是由“存在 \(\sigma \in \Gamma\) 使得 \(y = \sigma(x)\)”（X 的点 \(x\) 与 \(y\) 之间）所定义的等价关系。证明商空间 X/R 是豪斯多夫（相应地，正则）的。
\end{enumerate}

* 10) 设 X 是 R 的这样的子空间：它是点 \(\mathbf{I} / n\) 所成之集的余集，其中 \(n\) 是异于 \(\mathbf{o}, \mathbf{i}\) 和 \(-\mathbf{i}\) 的有理整数。设 S 是 X 上的等价关系，其等价类为 (i) 集合 \(\{\mathbf{o}\}\)；(ii) 全体非零整数所成之集；(iii) 所有形如 \(\{x, \mathbf{i} / x\}\) 的集合，其中 \(x \in \mathbf{X}\) 且 \(\mathbf{o} < |x| < \mathbf{i}\)。

\begin{enumerate}
\def\labelenumi{\alph{enumi})}
\item
  证明 X 上的关系 S 既非开的也非闭的。
\item
  证明 S 的图像在 \(X \times X\) 中闭，但商空间 X/S 不是豪斯多夫的。*
\end{enumerate}

¶ * 11) 设 X 是拓扑空间。用 Z 表示积空间 \(R^{X}\)，并对每个 \(x \in X\)，设 \(Z_{x}\) 为 Z 中由所有满足如下条件的 \(u \in R^{X}\) 组成的子集：\(u(x) \in Q\)，且 \(u(y) \notin Q\) 对所有异于 x 的 \(y \in X\) 成立。设 Y 是积 \(X \times Z\) 的子空间，它是诸集 \(\{x\} \times Z_x\)（当 \(x\) 取遍 X 时）的并。证明 Y 是豪斯多夫的，且 X 同胚于 Y 的一个商空间。*

\begin{enumerate}
\def\labelenumi{\arabic{enumi})}
\setcounter{enumi}{11}
\tightlist
\item
  设 X 是至少有两点的拓扑空间。
\end{enumerate}

\begin{enumerate}
\def\labelenumi{\alph{enumi})}
\item
  证明拓扑 \(\mathfrak{C}_{\Omega}\) 和 \(\mathfrak{C}_{\Phi}\)（在 \(\mathfrak{P}_0(\mathbf{X})\) 上；§ 2，习题 7）都不是可达空间的拓扑。
\item
  设 \(\mathcal{G}_{\Theta}\) 表示拓扑 \(\mathcal{G}_{\Omega}\) 与 \(\mathcal{G}_{\Phi}\) 的上确界。证明集 \(\mathfrak{F}(\mathbf{X})\)（\(\mathbf{X}\) 的非空闭子集所成之集）赋以由 \(\mathcal{G}_{\Theta}\) 诱导的拓扑后是 Kolmogoroff 空间；且若 \(\mathbf{X}\) 可达，则 \(\mathfrak{F}(\mathbf{X})\) 也可达。
\item
  在 X 可达的假设下，证明空间 \(\mathfrak{F}(\mathbf{X})\)（赋以由 \(C_{\Theta}\) 诱导的拓扑）是豪斯多夫的当且仅当 X 正则。
\end{enumerate}

¶ 13) 设 X 是一个集合，Γ 是 X 的一个置换群。a) 设 F 是 X 的子集所成的一个集。设 S 是 X 上使 F 中所有集合皆闭且每个 u ∈ Γ 都是同胚的所有拓扑组成的集。证明集 S 有一个最小元 C(Γ, F)。若 X 无限，证明存在 X 的子集所成的集 F，使得 C(Γ, F) 可达但非离散（§ 2，习题 8）。

\begin{enumerate}
\def\labelenumi{\alph{enumi})}
\setcounter{enumi}{1}
\tightlist
\item
  设 \(J(\Gamma)\) 是具有下述性质的子集 \(A\)（属于 \(X\)）所成的集：对每个 \(x \notin A\)，存在映射 \(f \in \Gamma\)，使得 \(f(x) \neq x\) 且保持 \(A\) 的每个元素不动。证明 \(X\) 上使每个映射 \(u \in \Gamma\) 都是 \(X\) 到自身的同胚的每个豪斯多夫拓扑，都比 \(\mathcal{C}_{J}(\Gamma) = \mathcal{C}(\Gamma, J(\Gamma))\) 更细。
\end{enumerate}

证明拓扑 \(\mathfrak{G}_{\mathrm{J}}(\Gamma)\) 是豪斯多夫的当且仅当群 \(\Gamma\) 满足下列条件：对任意两个不同的元素 \(x, y\)（属于 \(\mathbf{X}\)），存在有限个元素 \(u_{k} \in \Gamma\)，使得（记 \(\mathbf{F}(u_{k})\) 为 \(\mathbf{X}\) 中在 \(u_{k}\) 下不变的元素所成之集）诸 \(\mathbf{F}(u_{k})\) 覆盖 \(\mathbf{X}\)，且没有一个 \(\mathbf{F}(u_{k})\) 同时含 \(x\) 和 \(y\)。

\begin{enumerate}
\def\labelenumi{\alph{enumi})}
\setcounter{enumi}{2}
\item
  设 \(\mathcal{C}_{0}\) 表示有理直线 \(\mathbf{Q}\)（*相应地，实直线 \(\mathbf{R}_{*}\)*）的拓扑，并设 \(\Gamma\) 是 \(\mathbf{Q}\)（*相应地，\(\mathbf{R}_{*}\)*）到自身的所有同胚组成的群。证明拓扑 \(\mathcal{C}_{\mathbf{J}}(\Gamma)\) 与 \(\mathcal{C}_{0}\) 相同。
\item
  设 X 是有理直线的由点 0 和诸点 \(\iota/n\)（n 为整数 \(\geqslant\iota\)）组成的子空间，并设 \(\Gamma\) 是 X 的同胚群。证明拓扑 \(\mathcal{G}_{\mathrm{J}}(\Gamma)\) 不是豪斯多夫的，且 X 关于这个拓扑的同胚群就是 \(\Gamma\)。
\item
  设 X 是这样的拓扑空间：对 X 的每对不同的点 x, y，存在 X 到自身的同胚 u，使得 \(u(x) = x\) 且 \(u(y) \neq y\)（*例如空间 \(R^{n}_{*}\)*）。设 G 是 X 到自身的同胚群，并对每个 \(x \in X\)，设 \(S_{x}\) 是 G 中保持 x 不动的子群。证明 \(S_{x}\) 等于它在 G 中的正规化子，且所有子群
\end{enumerate}

\(S_{x}\) 之交仅由单位元组成；特别地，G 的中心仅由单位元组成。

\begin{enumerate}
\def\labelenumi{\arabic{enumi})}
\setcounter{enumi}{13}
\tightlist
\item
  证明公理 \((\mathrm{O}_{\mathrm{III}})\) 与下列每条公理等价：
\end{enumerate}

\((\mathrm{O}_{\mathrm{III}}^{\prime\prime})\) 若 F 是 X 的任一闭子集，则 F 的所有闭邻域之交等于 F。

\((\mathrm{O}_{\mathrm{III}}^{\prime\prime})\) 若 F 是 X 上收敛于点 a 的任一滤子，则以 F 中诸集的闭包为基的 X 上的滤子 F 也收敛于 a。

\begin{enumerate}
\def\labelenumi{\arabic{enumi})}
\setcounter{enumi}{14}
\tightlist
\item
  \begin{enumerate}
  \def\labelenumii{\alph{enumii})}
  \tightlist
  \item
    证明满足公理 (O \(_{III}\) ) 的每个 Kolmogoroff 空间都是豪斯多夫的（从而是正则的）。
  \end{enumerate}
\end{enumerate}

\begin{enumerate}
\def\labelenumi{\alph{enumi})}
\setcounter{enumi}{1}
\tightlist
\item
  在一个三元集上构造一个满足公理 \((\mathrm{O}_{\mathrm{III}})\) 的非豪斯多夫拓扑。
\end{enumerate}

\begin{enumerate}
\def\labelenumi{\arabic{enumi})}
\setcounter{enumi}{15}
\item
  设 \((\mathbf{X}_t)\) 是拓扑空间的一个无限族，并在积集合 \(\mathbf{X} = \prod_{i} \mathbf{X}_i\) 上赋以§ 4 习题 9 中定义的拓扑之一。证明若诸 \(\mathbf{X}_t\) 是 Kolmogoroff（相应地，可达、豪斯多夫、正则）空间，则 \(\mathbf{X}\) 亦然。
\item
  证明拓扑 \(\mathcal{C}_0(\mathbf{X})\)、\(\mathcal{C}_+(\mathbf{X})\) 和 \(\mathcal{C}_{-}(\mathbf{X})\)（§ 2，习题 5）在线性序集 \(\mathbf{X}\) 上都是正则的。
\item
  设 \(X_{1}\)，\(X_{2}\) 是两个集合，\(F_{i}\) 是 \(X_{i}\) 上的滤子（i = 1, 2），f 是 \(X_{1} \times X_{2}\) 到正则空间 Y 中的映射，使得对每个
\end{enumerate}

\[
x _ {1} \in \mathrm{X} _ {1} \text { there   exists } \lim _ {x _ {2}, \mathfrak {F} _ {2}} f (x _ {1}, x _ {2}) = g (x _ {1}).
\]

证明：点 \(a \in Y\) 是 \(g\) 关于滤子 \(\mathfrak{F}_1\) 的聚点，当且仅当任给一个邻域 \(V\)（\(a\) 在 \(Y\) 中的）和任一集 \(\mathbf{A}_1 \in \mathfrak{F}_1\)，存在 \(x_1 \in \mathbf{A}_1\) 及 \(\mathbf{A}_2 \in \mathfrak{F}_2\)，使得 \(f(\{x_1\} \times \mathbf{A}_2) \subset V\)。¶ 19) 设 \(X\) 是非正则的豪斯多夫空间（参见习题 20），\(a\) 是 \(X\) 的一点，使得存在邻域 \(U\)（\(a\) 的），它不含 \(a\) 的任何闭邻域。设 \(\mathfrak{F}\) 是 \(a\) 的邻域滤子，\(Y = \mathfrak{F} \cup \{\omega\}\) 是与 \(\mathfrak{F}\) 的截面滤子相伴随的拓扑空间。设 \(Z\) 是集合 \(Y \times X\)，赋以如下拓扑：对任何点 \((y, z) \neq (\omega, a)\)，\((y, z)\) 的邻域与积拓扑下的相同，而 \((\omega, a)\) 的邻域是含有形如

\[
(\{\omega \} \times \mathrm{V}) \cup ((\mathrm{Y} - \{\omega \}) \times \mathrm{X}),
\]

之集，其中 \(\mathbf{V}\) 是 \(a\) 在 \(\mathbf{X}\) 中的邻域。设 \(\mathbf{A}\) 是 \(\mathbf{Z}\) 的子集，它是诸集 \(\{\mathbf{V}\} \times \mathbf{V}\)（\(\mathbf{V} \in \mathfrak{F}\)）的并。设 \(f\) 是 A 到 X 中的映射 \((V, x) \rightarrow x\)。证明 f 在 \(\overline{A}\) 的每一点处关于 A 都有极限，但不存在 \(\overline{A}\) 到 X 中的连续映射延拓 f。

¶ 20) a) 拓扑空间 X 的开子集 U 称为正则的，如果 \(\mathbf{U} = \alpha(u)\)（§ 1，习题 3）；等价地，U 是正则的开集，如果 U 是某个闭集的内部。如果 X 的正则开子集构成 C 的一个基，则称 X 是半正则的（并称它的拓扑 C 是半正则的）。证明：若 C 满足公理 \((\mathrm{O}_{\mathrm{III}})\)，则 C 是半正则的。

\begin{enumerate}
\def\labelenumi{\alph{enumi})}
\setcounter{enumi}{1}
\item
  证明两个正则开集的交是正则开集。由此推出：关于 \(\mathcal{C}\) 的正则开集构成 X 上一个拓扑 \(\mathcal{C}^*\) 的基，该拓扑比 \(\mathcal{C}\) 粗且半正则；\(\mathcal{C}^*\) 称为与 \(\mathcal{C}\) 相伴随的半正则拓扑。证明：对 X 的每个在 \(\mathcal{C}\) 中开的子集 U，U 关于 \(\mathcal{C}^*\) 的闭包与它在 \(\mathcal{C}\) 中的闭包相同；且若 X 关于 \(\mathcal{C}\) 可达，则 X 关于 \(\mathcal{C}^*\) 的孤立点与关于 \(\mathcal{C}\) 的孤立点相同。证明 \(\mathcal{C}^*\) 是豪斯多夫的当且仅当 \(\mathcal{C}\) 是豪斯多夫的{[}利用§ 1 的习题 3 d){]}；最后证明：若 Y 是正则空间，则 X 到 Y 中的连续映射在拓扑 \(\mathcal{C}\) 与 \(\mathcal{C}^*\) 下是相同的。
\item
  设 X 是半正则空间，\(C_{0}\) 是它的拓扑。证明 X 上的拓扑 C 满足 \(C^{*} = C_{0}\)，当且仅当存在 X 的稠密子集（在拓扑 \(C_{0}\) 下）所成的一个集 M，使得 M 中诸集的每个有限交属于 M，且拓扑 C 由 M 与在拓扑 \(C_{0}\) 下为开的 X 的子集所成之集的并生成。{[}考虑（拓扑 C 下的）稠密开子集，并注意 C 中的每个开集是一个（拓扑 C 下的）稠密开集与一个（拓扑 \(C_{0}\) 下的）开集的交。{]}由此构造比正则拓扑更细的非正则豪斯多夫拓扑的例子（*甚至比可度量化紧空间的拓扑更细*）。
\end{enumerate}

\begin{enumerate}
\def\labelenumi{\arabic{enumi})}
\setcounter{enumi}{20}
\tightlist
\item
  设 X 是没有孤立点的豪斯多夫空间，\(C_{0}\) 是 X 的拓扑。设 M 表示 X 的这样的可数子集 A 的余集所成之集：\(\overline{A}\) 只有有限个非孤立点。证明：若 C 是 X 上由 M 与在 \(C_{0}\) 下为开的子集所成之集的并生成的拓扑，则 \(C^{*} = C_{0}\)（习题 20），且关于 C 收敛的每个序列只有有限个不同的项（这蕴含关于 C，X 的任何点都没有可数的基本邻域系，尽管 X 本身可以是可数的，从而 X 的每一点都是该点的可数个邻域之交）。
\end{enumerate}

¶ 22) a) 用习题 20 c) 的记号，证明拓扑的集 \(\mathrm{E}(\mathfrak{C}_{0})\)（由满足 \(C^{*}=C_{0}\) 的拓扑 C 组成）关于关系“C 比 \(C'\) 粗”是归纳的。若 \(C_{0}\) 是 X 上的半正则拓扑，则 \(\mathrm{E}(\mathfrak{C}_{0})\) 的一个极大元称为次极大拓扑，而赋以这种拓扑的集合 X 称为次极大空间。于是 \(\mathrm{E}(\mathfrak{C}_{0})\) 中的每个拓扑都比某个次极大拓扑粗。

\begin{enumerate}
\def\labelenumi{\alph{enumi})}
\setcounter{enumi}{1}
\item
  \(\mathfrak{C}\) 是 \(\mathbf{X}\) 上的次极大拓扑，当且仅当 \(\mathbf{X}\) 的每个在拓扑 \(\mathfrak{C}\) 下稠密的子集在 \(\mathfrak{C}\) 中都是开的。因此非空的次极大空间不是可解的（§ 2，习题 2）。
\item
  证明次极大空间的每个子空间是局部闭且次极大的（先考虑开子空间与闭子空间）。
\item
  证明与一个滤子相伴随的空间是次极大空间，这里该滤子比无限集的有限子集之余集所成的滤子更细。
\item
  若 M 是次极大空间 X 的子集，证明子空间 \(\overline{M} - \mathring{M} = \operatorname{Fr}(M)\) 是离散的。
\item
  反之，设 X 是这样的拓扑空间：它有一个稠密开子集 U，使得 U 是次极大的，且 X 的拓扑是 U-极大的（§ 3，习题 11）。那么 X 是次极大的。
\end{enumerate}

\begin{enumerate}
\def\labelenumi{\arabic{enumi})}
\setcounter{enumi}{22}
\item
  设 X 是拓扑空间，A 是 X 的稠密子集，使得 X 的拓扑是 A-极大的（§ 3，习题 11）。设 f 是 A 到拓扑空间 Y 中的连续映射，使得在 X --- A 的每一点处，f 关于 A 的聚点集非空（*例如当 X 非空且 Y 拟紧时即是这种情形*）。证明 f 可以连续延拓到整个 X 上。
\item
  在有理直线的区间 {]}0,1{[} 中，设 A 是形如 \(k/2^n\) 的有理数所成之集，B 是形如 \(k/3^n\) 的有理数所成之集（\(k\) 为整数）。设 X 是在 {]}0,1{[} 上赋以如下拓扑而得的空间：该拓扑由包含在 X 中的开区间 {]}α,β{[} 以及集合 A 和 B 生成。证明 X 是豪斯多夫且半正则的（习题 20），但不是正则的。
\end{enumerate}

¶ 25 a) 在集合 X 上，正则拓扑所成之集与没有孤立点的正则拓扑所成之集关于关系“C 比 \(C'\) 粗”都是归纳集。

\begin{enumerate}
\def\labelenumi{\alph{enumi})}
\setcounter{enumi}{1}
\item
  集合 X 上的拓扑 \(\mathcal{G}\) 称为超正则的，如果它在 X 上正则且无孤立点的拓扑所成的集中是极大的；因此对 X 上任何无孤立点的正则拓扑，都存在比它更细的超正则拓扑。空间称为超正则的，如果它的拓扑是超正则的。证明无孤立点的正则拓扑 \(\mathcal{G}\) 是超正则的，当且仅当只要 A 与 B 是 X 的两个互补的无孤立点子集，A 和 B 就都是开集。特别地，超正则空间不是可解的（§ 2，习题 11）。（为证条件是充分的，注意：若 \(\mathfrak{C}'\) 是比 \(\mathfrak{G}\) 更细的无孤立点正则拓扑，U 是关于 \(\mathfrak{C}'\) 的非空开集，\(\overline{\mathrm{U}}\) 是它在拓扑 \(\mathfrak{C}'\) 下的闭包，则 \(\overline{\mathrm{U}}\) 与 \(\mathrm{X} - \overline{\mathrm{U}}\) 关于拓扑 \(\mathfrak{C}'\) 都没有孤立点。）
\item
  在超正则空间中，每个无孤立点的集合的闭包是开的，且稠密集的内部稠密；因此两个稠密集之交稠密。由此推出：若 \(C_{0}\) 是超正则拓扑，则在所有满足 \(C^{*} = C_{0}\)（习题 20）的拓扑 C 之中，有一个（必然是次极大的）比所有其他拓扑都细。
\end{enumerate}

* 26) 设 \((\theta_{n})\) 是无理数序列，在区间 \(\mathbf{I} = [\mathbf{o},\mathbf{i}]\)（\(\mathbf{R}\) 中的）内稠密。设 \(\mathbf{I}_n\) 是把 \(\mathbf{I}\) 的所有点 \(0, \theta_1, \ldots, \theta_n\) 等同起来而得的商空间，\(\varphi_n\) 是典范映射 \(\mathbf{I} \to \mathbf{I}_n\)。设 \(\mathbf{X}\) 是含于 \(\mathbf{I}\) 的有理数所成之集；则 \(\varphi_n\) 在 \(\mathbf{X}\) 上的限制是 \(\mathbf{X}\) 到 \(\varphi_n(\mathbf{X})\) 上的双射。设 \(\mathcal{G}_n\) 是在 \(\varphi_n(\mathbf{X})\) 上由 \(\mathbf{I}_n\) 的拓扑诱导的拓扑经该双射的逆像。证明递减序列 \((\mathcal{G}_n)\)（\(\mathbf{X}\) 上的豪斯多夫拓扑序列）的下确界不是豪斯多夫拓扑。*

¶ 27) 设 X 是拓扑空间。证明存在 Kolmogoroff（相应地，可达、豪斯多夫、正则）空间 Y 和连续映射 \(\varphi: X \to Y\)，它具有下列泛性质：对 X 到 Kolmogoroff（相应地，可达、豪斯多夫、正则）空间 Z 中的每个连续映射 \(f: X \to Z\)，存在唯一的连续映射 \(g: Y \to Z\)，使得 \(f = g \circ \varphi\)（《集合论》，第四章，§ 3，第 1 目）。证明：

\begin{enumerate}
\def\labelenumi{\alph{enumi})}
\item
  对 Kolmogoroff 空间，可取 Y 为 X 关于等价关系 \(\{\bar{x}\}=\{\bar{y}\}\) 的商空间。
\item
  对可达空间，可取 Y 为 X 关于下述等价关系 R 的商空间：R 的图像是 X 上所有满足 \(\mathrm{G}(x)\) 对每个 \(x \in X\) 在 x 中闭的等价关系之图像 G 的交；于是 R 的图像含有 x 与 y 之间的关系 \(\{\bar{x}\} \cap \{\bar{y}\} \neq \emptyset\) 的图像。
\item
  对豪斯多夫空间与正则空间，试证《集合论》第四章，§ 3，第 2 目的条件 \((\mathrm{CU}_{\mathbf{I}})\) 至 \((\mathrm{CU}_{\mathbf{III}})\) 都满足（特别要利用习题 6 a)）。给出这样的例子：X 不正则，而 X 到相应的泛正则空间中的映射 \(\varphi: X \to Y\) 是双射{[}参见习题 20 b){]}。
\end{enumerate}

\begin{enumerate}
\def\labelenumi{\alph{enumi})}
\setcounter{enumi}{3}
\tightlist
\item
  若 X 是满足公理 \((\mathrm{O}_{\mathrm{III}})\) 的拓扑空间，证明与 X 对应的泛 Kolmogoroff 空间典范同胚于与 X 对应的泛正则空间。
\end{enumerate}

\begin{enumerate}
\def\labelenumi{\arabic{enumi})}
\setcounter{enumi}{27}
\tightlist
\item
  设 X 是非正则的豪斯多夫空间，F 是 X 的闭子集，\(a \notin F\) 是 X 的一点，使得 a 的每个邻域都与 F 的每个邻域相交。设 R 是把 F 的所有点等同起来而得的 X 上的等价关系。证明 R 是闭的，且 R 的图像在 \(X \times X\) 中闭，但 R 不是豪斯多夫的。
\end{enumerate}

¶ 29) a) 设 R 是非空拓扑空间 X 上的豪斯多夫等价关系，\(\mathfrak{F}(X)\) 是 X 的非空闭子集所成的空间，赋以由 \(\mathfrak{G}_{\Theta}\)（习题 12）诱导的拓扑。证明每个位于 X/R 的闭包中之 \(A \in \mathfrak{F}(X)\)（{[}把 X/R 看作 \(\mathfrak{F}(X)\) 的子集{]}）都含于 R 的某个等价类中（用反证法）。

\begin{enumerate}
\def\labelenumi{\alph{enumi})}
\setcounter{enumi}{1}
\tightlist
\item
  再设 X 正则且 R 是开的。证明 X/R 于是是 \(\mathfrak{F}(X)\) 关于拓扑 \(\mathfrak{G}_{\Theta}\) 的闭子集（利用§ 7 第 6 目的命题 11）。
\end{enumerate}

\BourbakiExerciseGroup

\begin{enumerate}
\def\labelenumi{\arabic{enumi})}
\item
  设 X 为拓扑空间，S 为 X 的拓扑的一个子基（§ 2，第 3 目）。证明：如果由属于 S 的集组成的 X 的每个开覆盖都含有 X 的一个有限覆盖，则 X 是拟紧的。（注意：若 X 上的超滤子 U 不收敛，则每个 \(x \in X\) 都属于一个不属于 U 的 S 中的集；并利用 § 6 命题 5 的推论。）
\item
  证明：良序集 X 赋以拓扑 \(\mathcal{C}_{-}(\mathbf{X})\)（§ 2，习题 5）后是局部紧的，且它是紧的当且仅当 X 有最大元。由此推出，每个集都可以赋以拓扑，使之成为一个紧空间。
\item
  \begin{enumerate}
  \def\labelenumii{\alph{enumii})}
  \tightlist
  \item
    设 X 为豪斯多夫空间，A、B 为 X 的两个不相交的紧子集。证明：存在 A 的一个邻域和 B 的一个邻域，二者没有公共点。
  \end{enumerate}
\end{enumerate}

\begin{enumerate}
\def\labelenumi{\alph{enumi})}
\setcounter{enumi}{1}
\tightlist
\item
  设 X 为正则空间，A 为 X 的紧子集，B 为 X 的不与 A 相交的闭子集。证明：存在 A 的一个邻域和 B 的一个邻域，二者没有公共点。
\end{enumerate}

\begin{enumerate}
\def\labelenumi{\arabic{enumi})}
\setcounter{enumi}{3}
\item
  证明：在 § 7 习题 6 所定义的空间 X（它是正则的（§ 8，习题 16））中，每个紧子集都是有限的。
\item
  *a) 证明：§ 8 习题 7 所定义的非豪斯多夫可达空间 Y = X/S 中，Y 的每个点都有紧邻域；1 与 -1 在 Y 中的像 α、β 有非闭的紧邻域；α 的一个紧邻域与 β 的一个紧邻域的交不是拟紧的，而这两个邻域的并不是紧的。*
\end{enumerate}

\begin{enumerate}
\def\labelenumi{\alph{enumi})}
\setcounter{enumi}{1}
\item
  设 X 是 N 与一个无限集 A 的和。对每个 \(x \in X\)，设 \(\mathfrak{S}(x)\) 是如下定义的滤子基：若 \(x \in N\)，则 \(\mathfrak{S}(x)\) 由单个集 \(\{x\}\) 组成；若 \(x \in A\)，且用 \(S_{x,n}\) 表示以 x 及整数 \(\geqslant n\) 为元素的集，则 \(\mathfrak{S}(x)\) 由诸集 \(S_{x,n}\)（n 取遍 N）组成。证明：X 上存在一个拓扑，使得 \(\mathfrak{S}(x)\) 对每个 \(x \in X\) 都是 x 的基本邻域系，且在此拓扑下 X 是可达的但不是豪斯多夫的。再证明 X 不是拟紧的，但 X 有稠密的紧子集。
\item
  证明：集 X 上的两个拓扑，若 X 在其每个之下都是可达的且拟紧的，则它们可以是不同但可比较的{[}参见§ 8，习题 7 及 5 b){]}。
\end{enumerate}

\begin{enumerate}
\def\labelenumi{\arabic{enumi})}
\setcounter{enumi}{5}
\item
  设 X、Y 为两个拓扑空间，A（相应地，B）为 X（相应地，Y）的拟紧子集。证明：若 U 是 \(A \times B\) 在 \(X \times Y\) 中的任一邻域，则存在 X 中 A 的邻域 V 与 Y 中 B 的邻域 W，使得 \(V \times W \subset U\)。
\item
  \begin{enumerate}
  \def\labelenumii{\alph{enumii})}
  \tightlist
  \item
    试给出拟紧空间的一个逆系统 \((\mathbf{X}_{\alpha}, f_{\alpha\beta})\)，它的逆极限是空集（注意：在第 1 目例 2 的空间中，每个子空间都是拟紧的）。
  \end{enumerate}
\end{enumerate}

\begin{enumerate}
\def\labelenumi{\alph{enumi})}
\setcounter{enumi}{1}
\tightlist
\item
  设 \(p\) 为素数。用 \(Z_{n}\) 表示有理整数集 \(\mathbf{Z}\)，赋以 \(\mathbf{Z}\) 到 \(\mathbf{Z} / p^{n}\mathbf{Z}\) 上的典范映射下离散拓扑的逆像。若 \(m \leqslant n\)，用 \(f_{mn}\) 表示恒等映射 \(Z_{n} \to Z_{m}\)，它是连续的。证明：诸空间 \(Z_{n}\) 都是拟紧的，但逆系统 \((Z_{n}, f_{mn})\) 的逆极限不是拟紧的（*）。
\end{enumerate}

\begin{enumerate}
\def\labelenumi{\arabic{enumi})}
\setcounter{enumi}{7}
\tightlist
\item
  设 X 为拓扑空间，\((\mathrm{G}_{\alpha})_{\alpha\in\mathbb{A}}\) 为 X 上等价关系图像的一个族，其指标集 A 是有向的；设 \(R_{\alpha}\) 为关系 \((x,y)\in\mathbf{G}_{\alpha}\)，\(\varphi_{\alpha}\) 为典范映射 \(X\to X/R_{\alpha}\)。设关系 \(R_{\beta}\) 蕴含 \(R_{\alpha}\)（只要 \(\alpha\leqslant\beta\)），并设 \(\varphi_{\alpha\beta}\) 为典范映射 \(X/R_{\beta}\to X/R_{\alpha}\)；于是 \((X/R_{\alpha},\varphi_{\alpha\beta})\) 是拓扑空间的一个逆系统。设 \(Y=\varprojlim X/R_{\alpha}\)。
\end{enumerate}

\begin{enumerate}
\def\labelenumi{\alph{enumi})}
\item
  设 \(g: \mathbf{X} \to \mathbf{Y}\) 为连续映射 \(\varprojlim \varphi_{\alpha}\)。证明：\(g(x)\) 在 \(\mathbf{Y}\) 中稠密，且对每个 \(x \in \mathbf{X}\)，\(\overrightarrow{g}(g(x))\) 是 \(x\) 关于其图像为 \(\bigcap_{n} G_{\alpha}\) 的等价关系的类。
\item
  设所有 \(R_{\alpha}\) 都是闭等价关系，且对每个 \(x \in X\) 与 x 的每个邻域 V，存在指标 \(\alpha\)，使得 x 关于 \(R_{\alpha}\) 的类含于 V。证明：在这些条件下，g 是 X 到 \(g(\mathbf{X})\) 上的开映射。
\item
  证明：若模 \(R_{\alpha}\) 的各类对每个 \(\alpha \in A\) 都是闭的且拟紧的，则 g 是满射。
\end{enumerate}

\begin{enumerate}
\def\labelenumi{\arabic{enumi})}
\setcounter{enumi}{8}
\item
  设 X 为可达空间，\(\mathcal{R}\) 为 X 的拟紧闭子集所成的一个集，它包含 X 的所有有限子集，且 \(\mathcal{R}\) 的任意两个子集之并属于 \(\mathcal{R}\)。设 \(X'\) 为一个集，它是 X 与由单点组成的集 \(\{\omega\}\) 的和。证明：\(X'\) 上存在一个拟紧的可达拓扑，它在 X 上诱导出给定的拓扑，且使得 \(X'\) 中 \(\mathcal{R}\) 的诸集之余集构成 \(\omega\) 的开邻域的一个基本系。试给出 \(X'\) 上与不同的集 \(\mathcal{R}\) 相对应的不同拓扑的例子。
\item
  设 X 为仿紧空间。证明：若 X 的每个开子空间都是仿紧的，则 X 的每个子空间都是仿紧的。由此推出：拓扑具有可数基的局部紧空间，其每个子空间都是仿紧的{[}参见第九章，§ 4，习题 20 a){]}。
\end{enumerate}

¶ 11) 设 \(\mathbf{X}_0\) 为具有最大元的不可数良序集；设 \(a\) 为 \(\mathbf{X}_0\) 的最小元，\(b\) 为元素 \(x \in \mathbf{X}_0\) 中使区间 \([a, x[\) 不可数的最小者。用 \(\mathbf{X}\) 表示区间 \([a, b[\)，赋以拓扑 \(\mathfrak{C}_{-}(\mathbf{X})\)。

\begin{enumerate}
\def\labelenumi{\alph{enumi})}
\item
  X 是局部紧的但不是紧的（习题 2），且 X 中每个序列都有聚点（注意 X 的每个可数子集都是上有界的）。
\item
  设 \(f\) 是 \(\mathbf{X}\) 到自身的映射，使得 \(f(x) < x\) 对一切充分大的 \(x\) 成立。证明：存在元素 \(c \in \mathbf{X}\)，具有如下性质：对每个 \(x \in \mathbf{X}\)，存在 \(y \geqslant x\) 使得 \(f(y) \leqslant c\)。{[}用反证法，构造点序列 \((z_n)\)（\(\mathbf{X}\) 的点），使得 \(z_{n+1}\) 是元素 \(z' \in \mathbf{X}\) 中满足 \(f(x) \geqslant z_n\)（对每个 \(x \geqslant z'\)）的最小者{]}。
\item
  由 b) 推出 X 不是仿紧的。
\end{enumerate}

\begin{enumerate}
\def\labelenumi{\arabic{enumi})}
\setcounter{enumi}{11}
\tightlist
\item
  设 X 为拓扑空间，\(\mathfrak{F}(X)\) 为 X 的非空闭子集所成之集，\(\mathfrak{K}(X)\) 为 X 的非空拟紧子集所成之集；
\end{enumerate}

给这些集中的每一个赋以 \(\mathfrak{C}_{\Omega}\)（\(\mathfrak{P}_{0}(\mathbf{X})\) 上的拓扑）所诱导的拓扑（§ 2，习题 7）。

\begin{enumerate}
\def\labelenumi{\alph{enumi})}
\tightlist
\item
  证明：若 \(\mathfrak{B} \subset \mathfrak{F}(\mathbf{X})\) 关于 \(\mathcal{C}_{\Omega}\) 是拟紧的，且 \(\mathbf{X}\) 是正则的，则诸集 \(\mathbf{M} \in \mathfrak{B}\) 的并在 \(\mathbf{X}\) 中是闭的。\\
\item
  证明：若 \(\mathfrak{B} \subset \mathfrak{K}(\mathbf{X})\) 关于 \(\mathcal{C}_{\Omega}\) 是拟紧的，则诸集 \(\mathbf{M} \in \mathfrak{B}\) 的并是拟紧的。
\end{enumerate}

¶ 13) 用习题 13 的记号，给 \(\mathfrak{F}(\mathbf{X})\) 和 \(\mathfrak{R}(\mathbf{X})\) 赋以 \(\mathfrak{C}_{\Theta}\)（\(\mathfrak{P}_{0}(\mathbf{X})\) 上的拓扑）所诱导的拓扑（§ 8，习题 12）。a) 证明：若 X 是拟紧的，则 \(\mathfrak{F}(\mathbf{X})\) 也是。（对 X 的每个开集 U，设 \(c(\mathrm{U})\){[}相应地 \(m(\mathrm{U})\){]}为 X 中所有满足 \(\mathrm{M} \subset \mathrm{U}\)（相应地 \(\mathrm{M} \cap \mathrm{U} \neq \emptyset\)）的闭子集 M 所成之集；诸 \(c(\mathrm{U})\) 与 \(m(\mathrm{U})\) 构成 \(\mathfrak{C}_{\Theta}\) 的一个子基。然后利用习题 1）。

\begin{enumerate}
\def\labelenumi{\alph{enumi})}
\setcounter{enumi}{1}
\item
  设 X 是可达的，并设 \(X'\) 是 X 在映射 \(x \to \{x\}\)（X 到 \(\mathfrak{F}(X)\) 中的映射）下的像。设 H 为 \(\mathfrak{F}(X)\) 的一个包含 \(X'\) 的子空间。证明：若 H 是拟紧的，则 X 也是。{[}若 \((\mathrm{U}_{\alpha})\) 是 X 的一个开覆盖，则注意：若用 \(B_{\alpha}\) 表示所有 \(M \in \mathfrak{F}(X)\) 中满足 \(M \cap U_{\alpha} \neq \emptyset\) 者所成之集，则 \(\{\mathfrak{B}_{\alpha}\}\) 是 \(\mathfrak{F}(X)]\) 的一个开覆盖。
\item
  证明：若 X 是局部紧的，则 \(\mathfrak{K}(\mathbf{X})\) 在 \(\mathfrak{F}(\mathbf{X})\) 中是开的（利用 § 7 命题 10）。
\item
  设 X 是可达的。证明：\(\mathcal{R}(X)\) 是豪斯多夫的（相应地，正则的、紧的、局部紧的）当且仅当 X 如此。{[}对前两个断言，利用 § 8 习题 12 与 § 9 习题 3；对后两个断言，利用 a)、b) 及习题 12{]}。
\end{enumerate}

¶ 14) 拓扑空间 X 称为 Lindelöf 空间，如果 X 的每个开覆盖都含有 X 的一个可数覆盖。每个具有可数基的空间都是 Lindelöf 空间，每个拟紧空间也是。

\begin{enumerate}
\def\labelenumi{\alph{enumi})}
\item
  Lindelöf 空间的每个闭子空间都是 Lindelöf 空间。为使 Lindelöf 空间的每个子空间都是 Lindelöf 空间，必须且只需每个开子空间都是 Lindelöf 的（参见习题 15）。
\item
  若 \(f\) 是 Lindelöf 空间 \(\mathbf{X}\) 到拓扑空间 \(\mathbf{X}'\) 中的任一连续映射，则子空间 \(f(\mathbf{X})\)（\(\mathbf{X}'\) 的）是 Lindelöf 的。
\item
  拓扑空间的 Lindelöf 子空间的每个可数并都是 Lindelöf 的。特别地，每个可数空间都是 Lindelöf 的。
\item
  Lindelöf 空间 X 若 X 中每个序列都有聚点，则是拟紧的{[}证明此时公理 (C'') 成立{]}。
\end{enumerate}

¶ 15) a) 设 X 为每个开子空间都是 Lindelöf 的拓扑空间。证明：X 满足 § 1 习题 7 中的条件 (D \(_{III}\) )，且若 X 是正则的，则 X 的每个闭子集都是开集的可数交{[}参见习题 23 c){]}。

\begin{enumerate}
\def\labelenumi{\alph{enumi})}
\setcounter{enumi}{1}
\item
  设 X 为每个开子空间都是仿紧的拓扑空间。证明：X 的每个开子空间都是 Lindelöf 的，当且仅当 X 满足 §1 习题 7 的条件 \((D_{\mathrm{III}})\)。
\item
  设 X 为紧空间。证明：X 的每个开子空间都是 Lindelöf 的，当且仅当 (i) X 满足 §1 习题 7 的条件 \((\mathbf{D}_{\mathbf{III}})\)，且 (ii) X 的每个闭子集都是开集的可数交（利用 b) 与第 10 目定理 5）。
\end{enumerate}

¶ 16) 设 X 为拓扑空间，A 为 X 的子集。点 \(a \in X\) 称为 A 的凝聚点，如果 \(a\) 的每个邻域都含有 A 的不可数无穷多个点。

\begin{enumerate}
\def\labelenumi{\alph{enumi})}
\item
  X 的子集 A 的凝聚点所成之集是闭的。试给出一个例子，其中 X 只有一个凝聚点（参见第 8 目定理 4）。
\item
  设 X 是每个开子空间都是 Lindelöf 的可达空间（习题 15）。证明：对 X 的每个子集 A，A 的凝聚点所成之集 B 是完备集，且 A ∩ GB 是可数的。
\end{enumerate}

¶ 17) 在拓扑空间 X 中，点 \(a\) 称为与子集 A 全然附着，如果对 \(a\) 的每个邻域 U 都有 card (A) = card (A ∩ U)。证明：X 是拟紧的，当且仅当对 X 的每个无限子集 A，存在 X 的一点与 A 全然附着。（若 X 不是拟紧的，证明存在初始序数 \(\omega_{\alpha}\)（《集合论》第三章，§ 6，习题 10）以及 X 的开覆盖 (U \(_{\xi}\) )，其中 \(\xi\) 取遍序数集 \(<\omega_{\alpha}\)，使得对每个指标 \(\xi<\omega_{\alpha}\)，集 A \(_{\xi}\)（诸 U \(_{\eta}\)（\(\eta<\xi\)）之并在 X 中的余集）的势等于 N \(_{\alpha}\)；由此推出，可用超限归纳法定义 X 的点的一个族 ( \(x_{\xi}\) ) ( \(\xi<\omega_{\alpha}\) )，使得 \(x_{\xi}\in\mathrm{A}_{\xi}\) 对每个 \(\xi\) 成立，且 \(x_{\xi}\neq x_{\eta}\) 每当 \(<\eta\xi\)。）

¶ 18) 豪斯多夫空间 X 称为绝对闭的，如果它具有下述性质：对 X 到豪斯多夫空间 X' 的子空间上的每个同胚 f，f(X) 在 X' 中是闭的。每个紧空间都是绝对闭的。

\begin{enumerate}
\def\labelenumi{\alph{enumi})}
\tightlist
\item
  证明下列条件等价：
\end{enumerate}

(AF) 空间 X 是绝对闭的。

(AF') X 上由开集组成的每个滤子基至少有一个聚点。

(AF'') X 中交为空的闭子集族含有这样的有限子族：该子族中诸子集的内部之交为空。

(AF''') X 的每个开覆盖都含有诸闭包能覆盖 X 的有限子族。

\begin{enumerate}
\def\labelenumi{\alph{enumi})}
\setcounter{enumi}{1}
\item
  设 X 为豪斯多夫空间，U 为 X 的稠密开子集。设 U 上由开集组成的每个滤子基在 X 中至少有一个聚点。证明：在这些条件下 X 是绝对闭的。特别地，若 A 是绝对闭空间 X 的开子集，则 \(\overline{A}\) 是绝对闭的。
\item
  设 X 为绝对闭空间，f 为 X 到豪斯多夫空间 \(X'\) 中的连续映射。证明：\(f(\mathbf{X})\) 是 \(X'\) 的绝对闭子空间。
\item
  证明：两个绝对闭空间的积是绝对闭的（按第 10 目命题 17 的方式论证）。
\end{enumerate}

¶ 19) 豪斯多夫空间 X 称为极小的，如果它的拓扑 C 满足：X 上每个豪斯多夫且比 C 更粗的拓扑必与 C 相合；此时 C 必是半正则的（§ 8，习题 20）。证明：豪斯多夫空间 X 是极小的，当且仅当 X 上每个由开集组成且至多有一个聚点的滤子基都是收敛的。这相当于说：X 是绝对闭的（习题 18），且 X 上每个由开集组成且有单个聚点的滤子基都收敛于该点。

¶ 20) 拓扑空间 X 称为完全豪斯多夫的（且 X 的拓扑称为完全豪斯多夫的），如果 X 的每对不同的点都有不相交的闭邻域。

\begin{enumerate}
\def\labelenumi{\alph{enumi})}
\item
  每个正则空间都是完全豪斯多夫的。完全豪斯多夫空间的每个子空间都是完全豪斯多夫的。完全豪斯多夫空间的每个积都是完全豪斯多夫的。比完全豪斯多夫拓扑更细的每个拓扑都是完全豪斯多夫的。
\item
  每个极小的完全豪斯多夫空间 X（习题 19）都是紧的。（用反证法：考虑 X 上的超滤子 U 以及由属于 U 的开集所组成的滤子基）。
\end{enumerate}

\(^{*}c)\) 设 \(\theta\) 为无理数 \textgreater0。对每个点 \((x,y)\in\mathbf{Q}^{2}\) 及每个 \(\alpha>0\)，设 \(\mathrm{B}_{\alpha}(x,y)\) 为由 \((x,y)\) 与点 \((z,0)\in\mathbf{Q}^{2}\)（满足 \(|z-(x+\theta y)|<\alpha\) 或 \(|z-(x-\theta y)|<\alpha\)）所组成的集。设 \(\mathfrak{B}(x,y)\) 为诸 \(\mathrm{B}_{\alpha}(x,y)\)（\(\alpha\) 取遍实数集 \textgreater0）所成之集。证明：存在豪斯多夫拓扑 \(\mathcal{C}\)（在 \(\mathbf{Q}^{2}\) 上），使得对每个 \((x,y)\)，\(\mathfrak{B}(x,y)\) 都是 \((x,y)\) 关于 \(\mathcal{C}\) 的基本邻域系。证明：对 \(\mathbf{Q}^{2}\) 的任意两点 a、b，在拓扑 \(\mathcal{C}.*\) 下，a 的每个闭邻域都与 b 的每个闭邻域相交。

¶ 21) a) 设 X 为豪斯多夫空间，\(\mathfrak{G}\) 为 X 的拓扑。证明：X 是绝对闭的（习题 18）当且仅当半正则拓扑 \(\mathfrak{G}^*\)（与 \(\mathfrak{G}\) 相伴随）（ \(\S\) 8，习题 20）是某个极小空间的拓扑（习题 19）。

\begin{enumerate}
\def\labelenumi{\alph{enumi})}
\setcounter{enumi}{1}
\tightlist
\item
  若 \(\mathfrak{C}\) 是完全豪斯多夫的（习题 20），则 \(\mathfrak{C}_{*}\) 也是（利用 § 8 习题 20）。由此推出，绝对闭的正则空间是紧的。
\end{enumerate}

¶ 22) a) 设 X 为豪斯多夫空间，A 为 X 的稠密子集。证明：若 A 上的每个滤子都在 X 中至少有一个聚点，则 X 是绝对闭的{[}验证习题 18 的公理 (AF'){]}。

\begin{enumerate}
\def\labelenumi{\alph{enumi})}
\setcounter{enumi}{1}
\tightlist
\item
  设 \(X_{0}\) 为具有非开稠密子集 A 的紧空间，\(C_{0}\) 为 \(X_{0}\) 的拓扑。设 C 为 \(X_{0}\) 上由 A 与 \(X_{0}\) 的在 \(C_{0}\) 中为开的子集所生成的拓扑，于是 \(C^{*} = C_{0}\)（§ 8，习题 20）。若用 X 表示把集 \(X_{0}\) 赋以拓扑 C（习题 21）所得的非紧绝对闭空间，证明 A 上的每个滤子在 X 中都有聚点。
\end{enumerate}

¶ 23) a) 设 X 为可数的绝对闭空间。证明：X 的孤立点所成之集在 X 中稠密。{[}用反证法：把 X 的点排成一个序列，并用归纳法构造 X 的可数开覆盖 (U \(_{n}\) )，使得诸 \(\overline{U}_{n}\) 的任何有限并不覆盖 X{]}。

*b) 试给出一个不是 Lindelöf 空间的完全豪斯多夫绝对闭拓扑空间的例子{[}在积 \([0,1]\times[0,1]\) 中，考虑诸集 \(A\times\{o\}\) 的余集，其中 A 取遍 \([0,1]\) 的所有子集，并应用 § 8 习题 20 c) 的一般方法{]}。这个空间满足 § 1 习题 1 的条件 \((D_{II})\)，但不满足条件 \((D_{III})\)。

\begin{enumerate}
\def\labelenumi{\alph{enumi})}
\setcounter{enumi}{2}
\tightlist
\item
  设 \(X_{0}\) 为没有孤立点的紧空间，M 为 \(X_{0}\) 的可数子集的余集所成之集{[}于是 M 中的集都在 \(X_{0}\) 中稠密，参见 a){]}。设 C 为由 \(X_{0}\) 的开集与 M 中的集所生成的拓扑。证明：把 \(X_{0}\) 赋以拓扑 C 所得的空间 X 是绝对闭的且完全豪斯多夫的，并且是一个不满足 §1 习题 7 条件 ( \(D_{II}\) ) 的 Lindelöf 空间；再证明 X 的任何点都没有可数的基本邻域系，且 X 中任何各项皆不同的序列都没有聚点。若 \(X_{0}\) 的每个开子空间都是 Lindelöf 的{[}参见习题 16 c){]}，证明 X 亦然，因而特别地 X 满足 ( \(D_{III}\) ){[}习题 16 a){]}。当 \(X_{0} = [0,1]\) 时即属此情形。但证明在此情形下，存在 X 的可数闭子集，它们不是 X 的开集的可数交（参见第九章，§5，第 3 目）。*
\end{enumerate}

¶ 24) a) 设 X 为豪斯多夫空间，Γ 为 X 上这样的滤子全体所成之集：它们具有由开集组成的基，且在 X 中没有聚点。证明：集 Γ 若非空（即若 X 不是绝对闭的），则是归纳的。

\begin{enumerate}
\def\labelenumi{\alph{enumi})}
\setcounter{enumi}{1}
\tightlist
\item
  设 \(\Omega\) 为 \(\Gamma\) 中极大滤子所成之集，并设 \(X'\) 为 \(X\) 与 \(\Omega\) 的和所成之集。对每个 \(x \in X\)，用 \(\mathfrak{B}(x)\) 表示 \(x\) 在 \(X\) 中的邻域所成之集；对每个 \(\omega \in \Omega\)，用 \(\mathfrak{B}(\omega)\) 表示 \(X'\) 的形如 \(\{\omega\} \cup M\) 的子集所成之集，其中 \(M\) 属于滤子 \(\omega\)。于是 \(X'\) 上存在一个拓扑，使得 \(\mathfrak{B}(x')\) 是 \(x'\)（对每个 \(x' \in X'\)）的基本邻域系。证明：在此拓扑下，\(X'\) 是绝对闭的，且具有下述泛性质：对任何开连续映射 \(f: X \to Y\)（\(X\) 到绝对闭空间 \(Y\) 中的映射），存在唯一的连续映射 \(g: X' \to Y\) 延拓 \(f\)（考察诸滤子在 \(f\) 下的像，这些滤子属于 \(\Omega\)，并注意在绝对闭空间中，在具有由开集组成的基的滤子集中为极大的滤子必收敛）。特别地，\(X'\) 的拓扑是 X-极大的（ \(\S 3\) ，习题 11）。由此推出，存在拓扑为拟极大的绝对闭空间（ \(\S 2\) ，习题 6）。
\end{enumerate}

¶ 25) a) 设 X 为可达空间，\(\Phi\) 为 X 上具有由闭集组成的基的滤子所成之集。证明：集 \(\Phi\) 是归纳的，并设 \(X''\) 为 \(\Phi\) 中极大滤子所成之集。

\begin{enumerate}
\def\labelenumi{\alph{enumi})}
\setcounter{enumi}{1}
\tightlist
\item
  对 X 的每个开子集 U，设 \(U^{*}\) 为滤子 \(F \in X''\)（\(U \in F\) 者）所成之集。若 U、V 为 X 的两个开子集，则
\end{enumerate}

\[
(\mathrm{U} \cap \mathrm{V}) ^ {*} = \mathrm{U} ^ {*} \cap \mathrm{V} ^ {*} \quad \text { and } \quad (\mathrm{U} \cup \mathrm{V}) ^ {*} = \mathrm{U} ^ {*} \cup \mathrm{V} ^ {*}
\]

（注意：若 \(\mathfrak{F} \in \mathbf{X}''\) 满足 \(\mathbf{U} \notin \mathfrak{F}\)，则 \(\mathbf{X} = \mathbf{U} \in \mathfrak{F}\)）。证明：诸集 \(\mathbf{U}^*\) 构成 \(\mathbf{X}''\) 上一个拓扑的基，且在此拓扑下 \(\mathbf{X}''\) 是拟紧的{[}利用上述注意，验证公理 \((\mathbf{C}''')\){]}并且是可达的{[}若 \(\mathfrak{F}, \mathfrak{G}\) 是 \(\mathbf{X}''\) 的两个不同元素，证明存在两个不相交的闭子集 \(\mathbf{A}\) 与 \(\mathbf{B}\)（\(\mathbf{X}\) 中的），使得 \(\mathbf{A} \in \mathfrak{F}\) 且 \(\mathbf{B} \in \mathfrak{G}\){]}。

\begin{enumerate}
\def\labelenumi{\alph{enumi})}
\setcounter{enumi}{2}
\item
  对每个 \(x \in \mathbf{X}\)，设 \(\varphi(x)\) 为 \(\mathbf{X}\) 上以单个集 \(\{x\}\) 为基的超滤子。证明：\(\varphi\) 是 \(\mathbf{X}\) 到 \(\mathbf{X}''\) 的一个稠密子集上的同胚，且若 \(\mathbf{X}\) 是拟紧的（且可达的），则 \(\varphi(\mathbf{X}) = \mathbf{X}''\)。
\item
  证明：对每个连续映射 \(f\)（\(\mathbf{X}\) 到紧空间 \(\mathbf{Y}\) 中的映射），存在唯一的连续映射 \(g: \mathbf{X}'' \to \mathbf{Y}\)，使得 \(f = g \circ \varphi\)。
\item
  设 \(\tilde{\mathbf{X}}\) 为与 \(\mathbf{X}''\) 相伴随的泛豪斯多夫空间（§ 8，习题 27）；证明：典范映射 \(\psi: \mathbf{X}'' \to \tilde{\mathbf{X}}\) 是满射，且 \(\tilde{\mathbf{X}}\) 是紧的。于是，\(\mathbf{X}\) 到紧空间 \(\mathbf{Y}\) 中的每个连续映射都可以唯一地写成 \(f = h \circ \psi \circ \varphi\) 的形式，其中 \(h: \tilde{\mathbf{X}} \to \mathbf{Y}\) 是连续的。
\end{enumerate}

\begin{enumerate}
\def\labelenumi{\arabic{enumi})}
\setcounter{enumi}{25}
\tightlist
\item
  若 X 是离散空间，证明习题 25 所定义的空间 \(X''\) 与 \(\tilde{X}\) 相同，且若 C 是习题 24 所定义的空间 \(X'\) 的拓扑，则 \(X''\) 可等同于空间 \(X'\)（赋有与 C 相伴随的半正则拓扑 \(C^*\)）（§ 8，习题 20）。此外，集 \(X''\) 可典范地等同于 X 上的超滤子所成之集。紧空间 \(X''\) 称为 X 的超滤子空间。
\end{enumerate}

¶ 27) 证明：在紧空间 X 中，每个等势的{[}§ 2，习题 11 d){]}无限开子空间都是可解的。{[}利用第 2 目命题 1，证明存在 U 的拓扑的一个基 B，使得 card (B) = card (U)；然后利用 § 2 习题 11 c)。{]}由此推出，每个没有孤立点的局部紧空间都是可解的{[}利用 § 2 习题 11 d){]}；因而这样的空间不是次极大的。

\begin{enumerate}
\def\labelenumi{\arabic{enumi})}
\setcounter{enumi}{27}
\tightlist
\item
  \begin{enumerate}
  \def\labelenumii{\alph{enumii})}
  \tightlist
  \item
    设 K 为离散空间 \(\{0,1\}\)。若 X 是任一紧空间，设 F 为 K 到 X 的闭子集所成之集 \(\mathfrak{F}(X)\) 中的映射 f 全体所成之集，这些映射满足 \(\mathbf{X}=f(0)\cup f(1)\)。考察紧空间 \(Y=K^{F}\)；证明：对 Y 中的每个 \(y=(y_{f})_{f\in\mathbb{F}}\)，X 的闭集 \(f(y_{f})\) 之交至多由一个点组成。使此交非空的 \(y\in Y\) 所成之集 Z 是 Y 的闭子集。若 \(\varphi(y)\) 是诸 \(f(y_{f})\)（对 \(y\in Z\)）公共的唯一点，证明 \(\varphi\) 是 Z 到 X 上的连续满射。
  \end{enumerate}
\end{enumerate}

\begin{enumerate}
\def\labelenumi{\alph{enumi})}
\setcounter{enumi}{1}
\tightlist
\item
  试给出这样的紧空间 X 的例子：不存在形如 \(K^{A}\) 的紧空间到 X 上的连续满射{[}利用 § 4 习题 4 b) 与 § 9 习题 26{]}。
\end{enumerate}

§ 29) 拓扑空间 X 称为局部拟紧的，如果 X 的每个点都有拟紧邻域的一个基本系。a) 证明：在局部拟紧空间中，每个拟紧子集都有拟紧邻域的一个基本系。

\begin{enumerate}
\def\labelenumi{\alph{enumi})}
\setcounter{enumi}{1}
\tightlist
\item
  证明：局部拟紧空间的每个局部闭子空间都是局部拟紧的。
\end{enumerate}

\(^{*}c)\) 在闭欧几里得圆盘 B：\(||x|| \leqslant 1\)（实平面 \(R^{2}\) 中的）上，设 \(\mathcal{G}_{0}\) 为由 \(R^{2}\) 的拓扑诱导的拓扑，并设 \(\mathcal{G}\) 为如下定义的（更细的）拓扑：B 的点 \(x \neq (1,0)\) 的邻域滤子在 \(\mathcal{G}\) 下与在 \(\mathcal{G}_{0}\) 下相同，但对点 \((1,0)\)，拓扑 \(\mathcal{G}\) 中的基本邻域系由诸集

\[
\left\{\left(\mathrm{I}, \mathrm{o}\right) \right\} \cup \left(\mathbf {D} \cap \mathbf {B} _ {\mathbf {0}}\right),
\]

组成，其中 D 是以 (1,0) 为中心的任一开圆盘，而 \(B_{0}\) 是开圆盘 \(\|x\|<1\)。设 Y 是赋以拓扑 C 的集 B，并设 R 是 Y 上的等价关系，其等价类是集 S——由所有 \(x\neq(1,0)\)

（B 中使 \(\|x\| = 1\) 的点）组成——以及由 B - S 的单个点组成的 Y 的子集。证明：商空间 X = Y/R 是拟紧的但不是局部拟紧的。*

\BourbakiExerciseGroup

\begin{enumerate}
\def\labelenumi{\arabic{enumi})}
\tightlist
\item
  \begin{enumerate}
  \def\labelenumii{\alph{enumii})}
  \tightlist
  \item
    试给出连续映射 \(f: \mathbf{X} \to \mathbf{Y}\)、\(g: \mathbf{Y} \to \mathbf{Z}\) 的例子，使得 \(g \circ f\) 是逆紧的，\(f\) 不是满射，且 \(g\) 不是逆紧的。
  \end{enumerate}
\end{enumerate}

\begin{enumerate}
\def\labelenumi{\alph{enumi})}
\setcounter{enumi}{1}
\tightlist
\item
  试给出两个逆紧映射 \(f: \mathbf{X} \to \mathbf{Y}\) 与 \(g: \mathbf{X} \to \mathbf{Z}\)（其中 \(\mathbf{X}, \mathbf{Y}, \mathbf{Z}\) 是拟紧的）的例子，使得 \(x \to (f(x), g(x))\) 不是 \(\mathbf{X}\) 到 \(\mathbf{Y} \times \mathbf{Z}\) 中的逆紧映射。
\end{enumerate}

\begin{enumerate}
\def\labelenumi{\arabic{enumi})}
\setcounter{enumi}{1}
\tightlist
\item
  设 \(f: \mathbf{X} \to \mathbf{Y}\) 为逆紧映射，\(\mathbf{A}\) 为 \(\mathbf{X}\) 的子空间。证明：若 \(\mathbf{A} = \overline{f}^1(f(\mathbf{A})) \cap \overline{\mathbf{A}}\)，则限制 \(f|_{\mathbf{A}}\) 是 \(\mathbf{A}\) 到 \(f(\mathbf{A})\) 中的逆紧映射。这一条件是否必要？
\end{enumerate}

¶ 3) 设 \((\mathbf{X}_{\alpha}, f_{\alpha\beta})\) 为拟紧 Kolmogoroff 空间的一个逆系统，其中诸 \(f_{\alpha\beta}\) 是逆紧映射。证明：若诸 \(\mathbf{X}_{\alpha}\) 非空，则 \(\varprojlim \mathbf{X}_{\alpha}\) 非空（证明：在拟紧 Kolmogoroff 空间中，每个非空闭子集都含有由单个点组成的闭子集）。

\begin{enumerate}
\def\labelenumi{\arabic{enumi})}
\setcounter{enumi}{3}
\item
  试给出连续映射 \(f: \mathbf{X} \to \mathbf{Y}\)（其中 \(\mathbf{X}\) 与 \(\mathbf{Y}\) 是豪斯多夫空间）的一个例子，使得逆像 \(\overline{f}^1(\mathbf{K})\)（对每个紧子集 \(\mathbf{K}\)，\(\mathbf{Y}\) 中的）都是紧的，但该映射不是逆紧的（参见§ 9，习题 4）。
\item
  设 \(f: \mathbf{X} \to \mathbf{Y}\) 为逆紧映射。
\end{enumerate}

\begin{enumerate}
\def\labelenumi{\alph{enumi})}
\tightlist
\item
  证明：若 X 是正则的，则 \(f(\mathbf{X})\) 是 Y 的正则子空间。
\item
  证明：若 \(f(\mathbf{X})\) 是 Y 的 Lindelöf 子空间（§ 9，习题 14），则 X 是 Lindelöf 的（利用 § 5 第 4 目命题 10）。
\end{enumerate}

¶ 6) 设 X 是子空间 \(\mathbf{X}_i (i \in I)\) 的和所成的空间，\(f\) 是 X 到拓扑空间 Y 中的连续映射。对每个 \(i \in I\)，设 \(f_i: \mathbf{X}_i \to \mathbf{Y}\) 是 \(f\) 在 \(\mathbf{X}_i\) 上的限制。证明：\(f\) 是逆紧的，当且仅当每个 \(f_i\) 都是逆紧的，且族 \((f(\mathbf{X}_i))_{i \in I}\) 是局部有限的{[}利用定理 1 d){]}。

\begin{enumerate}
\def\labelenumi{\arabic{enumi})}
\setcounter{enumi}{6}
\item
  设 \(f: \mathbf{X} \to \mathbf{X}'\) 为逆紧映射，\(\mathbf{R}\)（相应地，\(\mathbf{R}'\)）为 \(\mathbf{X}\)（相应地，\(\mathbf{X}'\)）上与 \(f\) 相容的等价关系。设 \(\overline{f}: \mathbf{X} / \mathbf{R} \to \mathbf{X}' / \mathbf{R}'\) 为由 \(f\) 诱导的商空间之间的映射。设 (i) 在 \(f\) 下，关于 \(\mathbf{R}\) 饱和的每个集的像是关于 \(\mathbf{R}'\) 饱和的，且 (ii) \(\mathbf{R}'\) 是开的。证明：在这些条件下 \(\overline{f}\) 是逆紧的{[}利用定理 1 c){]}。
\item
  证明：在 § 5 第 2 目例 3 的假设下，典范映射 \(\mathbf{X} \to \mathbf{X} / \mathbf{R}\) 是逆紧的。
\end{enumerate}

¶ 9) a) 设 X、Y 为两个拓扑空间，\(f: \mathbf{X} \to \mathbf{Y}\) 为满射闭映射（但不必连续）。证明：若 \(\overline{f}^{1}(y)\) 对每个 \(y \in \mathbf{Y}\) 都是拟紧的，则 \(\overline{f}^{1}(\mathbf{K})\) 对 Y 的每个拟紧子集 K 都是拟紧的（或直接证明该结果，或利用 § 5 习题 7 与 § 9 习题 12）。

\begin{enumerate}
\def\labelenumi{\alph{enumi})}
\setcounter{enumi}{1}
\tightlist
\item
  再设 X 是正则的。证明：若 \(\overline{f}^{1}(y)\) 对每个 \(y \in Y\) 在 X 中是闭的，则 \(\overline{f}^{1}(K)\) 对 Y 的每个拟紧子集 K 在 X 中是闭的（方法同上）。又若 Y 是局部紧的，则 f 是连续的。
\end{enumerate}

\begin{enumerate}
\def\labelenumi{\arabic{enumi})}
\setcounter{enumi}{9}
\tightlist
\item
  设 X、Y 为两个拓扑空间。X 与 Y 之间的对应 (G, X, Y)（《集合论》第二章，§ 3，第 1 目）称为逆紧的，如果投影 \(\mathbf{pr}_2\) 在图像 G 上的限制是逆紧映射 \(\mathbf{G} \to \mathbf{Y}\)。
\end{enumerate}

\begin{enumerate}
\def\labelenumi{\alph{enumi})}
\item
  设 \(f: \mathbf{X} \to \mathbf{Y}\) 为映射。证明：\(f\) 是连续的，当且仅当对应 \(\overline{f}^1\)（\(\mathbf{Y}\) 与 \(\mathbf{X}\) 之间的）是逆紧的；且 \(f\) 是逆紧映射，当且仅当对应 \(\overline{f}^1\)（\(\mathbf{Y}\) 与 \(\mathbf{X}\) 之间的）与 \(f\)（\(\mathbf{X}\) 与 \(\mathbf{Y}\) 之间的）都是逆紧的。由此构造一个非连续映射 \(f\) 的例子，使对应 \(f\) 是逆紧的。
\item
  若 \((\mathbf{G}, \mathbf{X}, \mathbf{Y})\) 是逆紧对应，证明：只要 A 是 X 的闭子集，\(\mathbf{G}(\mathbf{A})\) 就是 Y 的闭子集。
\end{enumerate}

¶ 11) 设 (G, X, Y) 为逆紧对应（习题 10）。

\begin{enumerate}
\def\labelenumi{\alph{enumi})}
\tightlist
\item
  证明：若 \(\mathfrak{B}\) 是 \(X\) 上由闭集组成的滤子基，则
\end{enumerate}

\[
\bigcap_ {\mathbf {M} \in \mathfrak {B}} \mathrm{G} (\mathbf {M}) = \mathrm{G} \left(\bigcap_ {\mathbf {M} \in \mathfrak {B}} \mathbf {M}\right).
\]

\begin{enumerate}
\def\labelenumi{\alph{enumi})}
\setcounter{enumi}{1}
\tightlist
\item
  对每个 \(y \in \mathbf{Y}\) 及每个邻域 \(\mathbf{V}\)（\(\overline{\mathbf{G}}(y)\) 在 \(\mathbf{X}\) 中的），证明存在一个邻域 \(\mathbf{W}\)（\(y\) 在 \(\mathbf{Y}\) 中的），使得 \(\overline{\mathbf{G}}^{1}(\mathbf{W}) \subset \mathbf{V}\)（“根作为参数的函数的连续性”）（利用 § 9 习题 6）。
\end{enumerate}

¶ 12) a) 对应 (G, X, Y) 是逆紧的，当且仅当对每个拓扑空间 Z 及每个图像 E 在 Z × X 中为闭的对应 (E, Z, X)，图像 G \(\circ\) E 在 Z × Y 中是闭的。为证条件必要，注意 G \(\circ\) E 是 (E × Y) ∩ (Z × G) 在 \(\iota_{\mathbf{Z}} \times \mathrm{pr}_{2} : \mathbf{Z} \times \mathbf{G} \to \mathbf{Z} \times \mathbf{Y}\) 下的像。为证充分性，设 \(\delta_{\mathbf{Y}}\) 为对角映射 \(\mathbf{Y} \to \mathbf{Y} \times \mathbf{Y}\)；证明：对每个子集 \(\mathbf{F}\)（\((\mathbf{Z} \times \mathbf{Y}) \times \mathbf{X}\) 的），\(\mathbf{G} \circ \mathbf{F}\) 在映射

\[
\iota_ {\mathbf {Z}} \times \delta_ {\mathbf {Y}}: \mathbf {Z} \times \mathbf {Y} \rightarrow \mathbf {Z} \times \mathbf {Y} \times \mathbf {Y}
\]

下的逆像，是在映射 \((y, z) \to (z, y)\) 之下、集

\[
\operatorname{pr} _ {2 3} (\overline {{{\mathrm{F}}}} \cap (\mathrm{G} \times \mathrm{Z})),
\]

的像，其中 \(\mathbf{F}\) 是 \(\mathbf{F}\) 在映射 \((z, y, x) \to (x, y, z)\) 下的像。b) 由 a) 推出：两个逆紧对应的复合是逆紧对应。

\begin{enumerate}
\def\labelenumi{\alph{enumi})}
\setcounter{enumi}{2}
\tightlist
\item
  设 \((\mathbf{G},\mathbf{X},\mathbf{Y})\) 为逆紧对应。证明：若 \(\mathbf{X}\) 是豪斯多夫的，则 \(\mathbf{G}\) 在 \(\mathbf{X} \times \mathbf{Y}\) 中是闭的。
\end{enumerate}

¶ 13) 设 X 为局部紧空间，Y 为任一拓扑空间。a) 证明下列四个命题等价：

\(\alpha)\) (G, X, Y) 是 X 与 Y 之间的逆紧对应。

\(\beta)\) 对每个豪斯多夫空间 \(\mathbf{X}'\)（把 \(\mathbf{X}\) 作为子空间包含），\(\mathbf{G}\) 在 \(\mathbf{X}' \times \mathbf{Y}\) 中是闭的。

\(\gamma)\) G 在 \(\mathbf{X} \times \mathbf{Y}\) 中是闭的，且对每个 \(y \in \mathbf{Y}\)，存在邻域 \(\mathbf{V}\)（\(y\) 在 \(\mathbf{Y}\) 中的）与紧子集 \(\mathbf{K}\)（\(\mathbf{X}\) 的），使得 \((\mathbf{X} - \mathbf{K}) \times \mathbf{V}\) 不与 \(\mathbf{G}\) 相交。

δ) 存在紧空间 \(X'\)（把 X 作为子空间包含），使得 G 在 \(X' \times Y\) 中是闭的。

{[}为证 \(\alpha) \Longrightarrow \beta)\)，利用习题 12 \(a)\) 与 \(c)\)。为证 \(\delta) \Longrightarrow \alpha)\)，利用第 1 目命题 2 与第 2 目定理 1 的推论 5。{]}

\begin{enumerate}
\def\labelenumi{\alph{enumi})}
\setcounter{enumi}{1}
\tightlist
\item
  证明：若 \((\mathbf{G},\mathbf{X},\mathbf{Y})\) 是逆紧的，则：
\end{enumerate}

\(\zeta)\) 对每个紧子集 \(\mathbf{L}\)（\(\mathbf{Y}\) 的），\(\overline{\mathbf{G}}(\mathbf{L})\) 是紧集。

再证明：若 Y 是局部紧的，且 G 在 \(X \times Y\) 中是闭的并满足 \(\zeta\)，则 (G, X, Y) 是逆紧的。{[}利用 \(\alpha\) 与 \(\gamma\) 的等价性。{]}当 Y 是豪斯多夫但非局部紧的空间时，后一断言不再成立（习题 4）。

\begin{enumerate}
\def\labelenumi{\arabic{enumi})}
\setcounter{enumi}{13}
\item
  设 Y 为局部紧空间，X 为拓扑空间。证明：映射 \(f: X \rightarrow Y\) 是连续的，当且仅当对每个把 Y 作为子空间包含的紧空间 \(Y'\)，f 的图像在 \(X \times Y'\) 中是闭的。{[}利用习题 13 与 10 a){]}。
\item
  设 X 为豪斯多夫空间，A 为 X 的闭子集，f 为 X --- A 到局部紧空间 Y 中的逆紧映射，Y′ 为 Y 的单点紧化，\(\omega\) 为 Y′ 中的无穷远点。设 g 为 X 到 Y′ 中的映射，它在 X --- A 上与 f 相合，且把 A 的每个点映到 \(\omega\)。证明 g 是连续的。
\end{enumerate}

¶ 16) 设 X 为局部紧空间，R 为 X 上的等价关系，C 为 R 在 X × X 中的图像。证明：C 在 X × X 中是闭的，当且仅当对 X 的每个紧子集 K，K 关于 R 的饱和集在 X 中是闭的。（为证必要性，利用第 2 目定理 1 的推论 5）。

*17) 在局部紧空间 R 上，设 S 是把 Z 的所有点等同起来所得的等价关系。证明：S 是闭的，且 R/S 是豪斯多夫的但不是局部紧的。*

*18) 设 X 为 R 的局部紧子空间 {[}o, +∞{[}。设 S 为 X 上的等价关系，其等价类是诸集 \{x, 1/x\}（对 o \textless{} x ≤ 1）以及集 \{o\}。证明：S 是开的，S 的图像在 X × X 中是闭的，且 X/S 是紧的，但 S 不是闭的。*

¶ 19) 设 X 为局部紧 σ-紧空间，R 为 X 上其图像在 X × X 中为闭的等价关系。证明：X/R 是豪斯多夫的。{[}设 (Uₙ) 是 X 的由相对紧的开集组成的覆盖，使得 Uₙ ⊂ Uₙ₊₁。设 A、B 为两个关于 R 饱和的不相交闭集。递归地定义两个递增的、关于 R 饱和的闭集序列 (Vₙ)、(Wₙ)，使得 Vₙ ∩ Wₙ = ∅，且 Vₙ ∩ Uₙ（相应地，Wₙ ∩ Uₙ）是 A ∩ Uₙ（相应地，B ∩ Uₙ）在紧子空间 Uₙ 中的邻域。利用第 3 目命题 8 与习题 16，以及 § 9 第 2 目命题 2{]}。(*)

¶ 20) a) 设 X 为非空紧空间，R 为 X 上的豪斯多夫等价关系。证明：X/R 在空间 \(\mathfrak{P}_{0}(X)\) 中（赋有拓扑 \(\mathcal{G}_{\Theta}\)；§ 8，习题 12）是闭的，当且仅当关系 R 是开的。{[}利用 § 9 习题 13 与 § 3 习题 16 证明：若 X/R 在空间 \(\mathfrak{P}_{0}(X)\) 中（赋有拓扑 \(\mathcal{G}_{\Theta}\)）是闭的，则 X/R 上的商拓扑与由 \(\mathcal{G}_{\Theta}\) 诱导的拓扑相合，然后应用 § 5 习题 7。反之，若 R 是开的，利用 § 8 习题 29 证明 X/R 关于 \(\mathcal{G}_{\Theta}\) 是闭的{]}。

*b) 设 X 为实平面 \(\mathbb{R}^2\) 的局部紧子空间，它由诸点 \((\mathtt{I}/n, y)\)（其中 \(n\) 是整数 \(>0\) 且 \(-\mathtt{I} \leqslant y \leqslant \mathtt{I}\)）以及诸点 \((\mathtt{o}, y)\)（其中 \(-\mathtt{I} \leqslant y < 0\)，或 \(0 < y \leqslant \mathtt{I}\)，或 \(y = 2\)）组成。设 \(p\) 是 \(\mathsf{pr}_1\) 在 X 上的限制，并设 R 为与 \(p\) 相伴随的等价关系。证明：R 既非开亦非闭，但 X/R 是紧的，且在空间 \(\mathfrak{P}_0(X)\) 中（赋有拓扑 \(\mathcal{C}_{\Theta}\)）是闭的。* ¶ 21) (*) a) 设 X 为局部拟紧的可达空间（§ 9，习题 29）。设 \(\mathbf{X}_0\) 为与 X 有同一底集的离散空间，并设 \(\tilde{\mathbf{X}}_0\) 为 \(\mathbf{X}_0\) 的（紧）超滤子空间（把离散空间 \(\mathbf{X}_0\) 等同于 \(\tilde{\mathbf{X}}_0\) 的一个开子空间）（§ 9，习题 26）。在 \(\tilde{\mathbf{X}}_0\) 上，设 R {[}或 R(X){]} 为超滤子 U、B 之间的如下等价关系：“U 与 B 在 X 中的极限点所成之集相同”。证明：R 是豪斯多夫的{[}利用 § 9 习题 25 b)、§ 7 习题 7 b) 以及公理 (C') 中给出的 \(\tilde{\mathbf{X}}_0\) 的拓扑的描述{]}。紧空间 \(\mathbf{X}^c = \tilde{\mathbf{X}}_0 / \mathbf{R}\) 与 X 的本原子集所成之集（§ 7，习题 7）成一一对应。若 X 是豪斯多夫的（从而是局部紧的），则 \(\mathbf{X}^c\) 当 X 紧时可等同于 X，而当 X 非紧时可等同于 X 的单点紧化。

\begin{enumerate}
\def\labelenumi{\alph{enumi})}
\setcounter{enumi}{1}
\item
  设 Y 为 X 的闭子空间。超滤子空间 \(\tilde{Y}_{0}\) 可等同于 \(Y_{0}\) 在 \(\tilde{X}_{0}\) 中的闭包，而 \(Y^{c}\) 可等同于 \(\tilde{Y}_{0}\) 在 \(X^{c}\) 中的典范像。
\item
  证明：限制到 \(X_{0}\) 上的典范映射 \(\tilde{X}_{0} \rightarrow X^{c}\) 是 \(X_{0}\) 到子空间 \(X_{*}^{c}\)（\(X^{c}\) 的）上的双射；典范映射 \(f: X_{*}^{c} \rightarrow X\) 是连续双射；且 \(X_{*}^{c}\) 的任何紧子集在该映射下的像是 X 的闭紧子集。{[}为证 f 连续，利用 b)；还需注意：在 \(\tilde{X}_{0}\) 中，\(X_{*}^{c}\) 的一个点的逆像由 X 上具有同一个给定极限点的所有超滤子组成，且若 U 是 X 上的超滤子，则它作为 \(\tilde{X}_{0}\) 上的超滤子基在 \(\tilde{X}_{0}\) 中的极限点，就是把 U 视为 \(\tilde{X}_{0}\) 的点时所得的点{]}。
\end{enumerate}

*d) 设 \(\mathbf{Z}\) 为这样的集：它是可数个空间 \(\mathbf{Y}_n\)（每个都等同于 §8 习题 7 所定义的空间 \(\mathbf{Y}\)）与由单个点组成的集 \(\{\omega\}\) 的和。在 \(\mathbf{Z}\) 上如下定义一个拓扑：对点 \(z \in \mathbf{Y}_n\)，取它在空间 \(\mathbf{Y}_n\) 中的邻域滤子作为基本邻域系；对 \(\omega\)，取子集 \(\mathbf{V}_n\)（\(\mathbf{Y}\) 的）所成之集作为基本邻域系，其中 \(\mathbf{V}_n\) 是 \(\{\omega\}\) 与诸 \(\mathbf{Y}_m\)（\(m \geqslant n\) 者）的并。如此定义的空间 \(\mathbf{Z}\) 是可达的、拟紧的且局部拟紧的；但子空间 \(Z_*^c\)（\(Z^c\) 的）不是局部紧的。*

\begin{enumerate}
\def\labelenumi{\alph{enumi})}
\setcounter{enumi}{4}
\tightlist
\item
  设 Y 为另一个局部拟紧的可达空间，u 为 X 到 Y 中的一个连续映射，使得 Y 的每个拟紧子集在 u 下的逆像是拟紧的。把 u 视为 \(X_{*}^{c}\) 到 \(Y_{*}^{c}\) 中的映射，证明：u 是连续的，且可（唯一地）延拓为连续映射 \(\mathbf{X}^c\to \mathbf{Y}^c\)。{[}先把 \(u\) 延拓为 \(\tilde{\mathbf{X}}_0\) 到 \(\tilde{\mathbf{Y}}_0\) 中的一个连续映射，并证明该延拓与等价关系 \(\mathbf{R}(\mathbf{X}),\mathbf{R}(\mathbf{Y})\) 相容{]}。
\end{enumerate}

\(^{*}f)\) 设 X 为 § 9 习题 29 c) 所定义的拟紧、非局部拟紧的可达空间。证明：在相应的空间 \(\tilde{X}_{0}\) 中，等价关系 R(X) 不是豪斯多夫的。*

\BourbakiExerciseGroup

\begin{enumerate}
\def\labelenumi{\arabic{enumi})}
\item
  证明：§ 9 习题 20 c) 所定义的可数豪斯多夫空间是连通的。
\item
  \begin{enumerate}
  \def\labelenumii{\alph{enumii})}
  \tightlist
  \item
    设 X 为拓扑空间，\(\mathfrak{C}\) 为 X 的拓扑。证明：若 X 在半正则拓扑 \(\mathfrak{C}^*\)（\(\mathfrak{C}\) 的相伴随拓扑； \(\S 8\) ，习题 20）下是连通的，则 X 在原来的拓扑 \(\mathfrak{C}\) 下是连通的。由此推出：存在连通的次极大空间（ \(\S 8\) ，习题 22）。
  \end{enumerate}
\end{enumerate}

*b) 特别地，在实直线 \(\mathbf{R}\) 上构造一个拓扑，使得 \(\mathbf{R}\) 在其下是连通的，但 \(\mathbf{R}\) 的任何点都没有连通邻域的基本系{[}参见第四章，§ 2，习题 14 c){]}。*

\begin{enumerate}
\def\labelenumi{\arabic{enumi})}
\setcounter{enumi}{2}
\tightlist
\item
  设 A、B 为拓扑空间 X 的两个子集。
\end{enumerate}

\begin{enumerate}
\def\labelenumi{\alph{enumi})}
\item
  证明：若 A 与 B 在 X 中都是闭的，且 A ∪ B 与 A ∩ B 都是连通的，则 A 与 B 都是连通的。举例说明：当 A、B 中有一个不是闭集时，该命题不必再成立。
\item
  设 A 与 B 都是连通的，且 \(\overline{A} \cap B \neq \varnothing\)。证明：\(A \cup B\) 是连通的。
\end{enumerate}

\begin{enumerate}
\def\labelenumi{\arabic{enumi})}
\setcounter{enumi}{3}
\tightlist
\item
  设 X 为至少有两个点的连通空间。
\end{enumerate}

\begin{enumerate}
\def\labelenumi{\alph{enumi})}
\item
  设 A 为 X 的连通子集，B 为 \(\mathbb{C}\mathbf{A}\) 的一个在 \(\mathbb{C}\mathbf{A}\) 中既开又闭的子集。证明：\(\mathbf{A} \cup \mathbf{B}\) 是连通的。（把 §3 习题 5 应用于 \(\mathbf{Y} = \mathbf{A} \cup \mathbf{B}\) 与 \(\mathbf{Z} = \mathbb{C}\mathbf{A}\)。）
\item
  设 A 为 X 的连通子集，B 为集 CA 的一个分支。证明：CB 是连通的{[}利用 a){]}。
\item
  由 b) 推出：存在 X 的两个连通子集 M、N，它们都异于 X，且满足 \(M \cup N = X\) 与 \(M \cap N = \varnothing\)。
\end{enumerate}

*5) a) 在实平面 \(\mathbf{R}^2\) 中给出连通集的一个递减序列 \((\mathbf{A}_n)\) 的例子，使其交不是连通的（参见第二章，§ 4，习题 14）。

\begin{enumerate}
\def\labelenumi{\alph{enumi})}
\setcounter{enumi}{1}
\tightlist
\item
  在 \(\mathbf{R}^2\) 中，设 \(\mathbf{X}\) 为开半平面 \(y > 0\) 与 \(y < 0\) 以及诸点 \((n, 0)\)（其中 \(n\) 是整数 \(> 0\)）的并所成之集。
\end{enumerate}

X 在拓扑 \(\mathfrak{C}_0\)（由 \(\mathbb{R}^2\) 的拓扑诱导的）下是连通的。定义 X 上的一个拓扑序列 \((\mathfrak{C}_n)\)，使得 \(\mathfrak{C}_{n+1}\) 比 \(\mathfrak{C}_n\) 细，且 X 在每个 \(\mathfrak{C}_n\) 下都连通，但在序列 \((\mathfrak{C}_n)\) 的上确界拓扑下不连通{[}从 \(\mathfrak{C}_n\) 到 \(\mathfrak{C}_{n+1}\) 的过渡，是通过引入点 \((n+1,0)]\) 的新邻域实现的。*

\begin{enumerate}
\def\labelenumi{\arabic{enumi})}
\setcounter{enumi}{5}
\item
  设 X、Y 为两个连通空间，A（相应地，B）为 X（相应地，Y）的异于 X（相应地，Y）的子集。证明：在积空间 \(X \times Y\) 中，\(A \times B\) 的余集是连通的。
\item
  设 X、Y 为两个连通空间，f 为 \(X \times Y\) 到拓扑空间 Z 中的一个映射，使得所有偏映射 \(f(0, y) : X \to Z\)、\(f(x, 0) : Y \to Z (x \in Y, y \in Y)\) 都连续。证明：\(f(X \times Y)\) 是连通的。
\item
  在 § 4 习题 9 给出的积集 \(\mathbf{X} = \prod_{i\in I}\mathbf{X}_i\) 上的拓扑的例子中，设 I 是无限的且基数 t 是不可数的。证明：若诸 \(\mathbf{X}_i\) 都是连通的、正则的，且每个都至少有两个不同的点，则 X 的点 \(a = (a_i)\) 的连通分支是由这样的 \(x = (x_i)\) 所成之集：\(x_{i} = a_{i}\) 除去对有限多个指标 \(i\in I\) 外成立。（化归到 I = N 的情形；考察 X 的两个点 \(b = (b_n)\)、\(c = (c_n)\)，使得 \(b_{n}\neq c_{n}\) 对所有 n 成立。~对每个 n，设 \((\mathrm{V}_{nm})_{m\geqslant 0}\) 是 \(b_{n}\) 在 \(\mathbf{X}_n\) 中的开邻域序列，这些邻域不含 \(c_{n}\)，且使得 \(\overline{\mathbf{V}}_{n,m + 1}\subset \mathbf{V}_{nm}\)。设 Z 为 X 的具有下述性质的点 \((x_{n})\) 所成之集：存在整数 \(k > 0\)，使得 \(x_{n}\in \mathrm{V}_{n,n - k}\) 对所有满足 \(n\geqslant k\) 的 n 成立。证明：Z 在 X 中既开又闭。）
\end{enumerate}

*9) 证明：在 § 10 习题 20 所定义的局部紧空间 \(\mathbf{X}\) 中，存在点 \(x\)，其在 \(\mathbf{X}\) 中的分支不同于 \(\mathbf{X}\) 的既开又闭且含 \(x_*\) 的子集之交。

\begin{enumerate}
\def\labelenumi{\arabic{enumi})}
\setcounter{enumi}{9}
\item
  证明：线性有序集 X 赋以拓扑 \(\mathcal{G}_{+}(X)\) 或拓扑 \(\mathcal{G}_{-}(X)\)（§ 2，习题 5）后是全不连通的。
\item
  设 X 为拓扑空间，R 为 X 上的等价关系，使得每个 mod R 等价类都包含在 X 的一个分支中。证明：X/R 的诸分支是 X 的诸分支的典范像（参见第三章，§ 2，习题 17）。
\item
  设 X 为拓扑空间。证明下列三个条件等价：\(\alpha\) ) X 的诸分支是开集；\(\beta\) ) 每个 \(x \in X\) 都有一个邻域，含于每个既开又闭且含 \(x; \gamma\) ) 对每个 \(x \in X\) ，既开又闭且含 x 的诸集之交是一个开集。
\end{enumerate}

*13) 设 \(\theta\) 为一无理数。设 \(f\) 为实直线 \(\mathbf{R}\) 到环面 \(\mathbf{T}^2\) 中的映射，像 \(f(t)\)（每个 \(t \in \mathbf{R}\) 的）是 \(\mathbf{T}^2\) 中 \((t, \theta t) \in \mathbf{R}^2\) 的典范像。\(f\) 是单射且连续的。证明：在 \(\mathbf{R}\) 上赋以 \(f\) 下 \(\mathbf{T}^2\) 的拓扑的逆像拓扑时，\(\mathbf{R}\) 是连通的，但 \(\mathbf{R}\) 的任何点都没有连通邻域的基本系。*

\begin{enumerate}
\def\labelenumi{\arabic{enumi})}
\setcounter{enumi}{13}
\tightlist
\item
  \begin{enumerate}
  \def\labelenumii{\alph{enumii})}
  \tightlist
  \item
    设 X 为局部紧且局部连通的空间，A 为 X 的闭子集。设 A* 为 A 与 CA 的相对紧的诸分支之并。若 A* = A，则称 A 为满的。证明：A* 是闭的且满的。
  \end{enumerate}
\end{enumerate}

\begin{enumerate}
\def\labelenumi{\alph{enumi})}
\setcounter{enumi}{1}
\item
  再设 X 是连通的。证明：若 A 是 X 的紧子集，则 A* 是紧的（考察 A 的一个紧邻域 V 以及与 V 的边缘相交的 CA 的分支）。
\item
  再设 X 是连通且 \(\sigma\) -紧的。证明：存在满的连通紧集的一个递增序列 \((\mathbf{K}_{n})_{n\geqslant0}\)，其并为 X，且使得 \(K_{n}\subset\mathring{K}_{n+1}\) 对所有 n 成立。
\end{enumerate}

*d) 证明：若只设 X 是连通且局部紧的，或 X 是豪斯多夫、连通且局部连通的，则 b) 的结论不再成立。{[}考察 \(\mathbf{R}^2\) 的由这样的点 \((x, y)\) 组成的子空间：或者 \(y = 0\) 且 \(0 \leqslant x \leqslant 1\)，或者 \(y > 0\) 且 \(x = 1/n\)（其中 \(n\) 是整数 \(> 0\)）。*

¶ 15) 设 X 为连通的、局部连通的空间。

\begin{enumerate}
\def\labelenumi{\alph{enumi})}
\item
  设 M、N 为 X 的两个不相交的非空闭子集。证明：\(\mathbb{C}(M \cup N)\) 有一个分支，其闭包与 M 和 N 都相交。{[}若 C 是 \(\mathbb{C}(M \cup N)\) 的一个其闭包不与 M 相交的分支，注意 C 在 CN 中既开又闭{]}。
\item
  若 A 是 X 的异于 X 的闭子集，由 a) 推出：A 的每个分支都与 \(\mathbb{C}\mathbb{A}\) 的闭包相交。
\end{enumerate}

*16) 设 X 为实平面 \(\mathbf{R}^2\) 的子空间，它是直线 \(y = 0\)、以 (0, 1) 与 \((n, 1/(1 + n))\) 为端点的闭线段、以及半直线 \(x \leqslant n\)，\(y = 1/(1 + n)\)（ \(n\) 为整数 \(>0\)）的并。证明：X 是连通空间，但 X 中 (0, 1) 的余集有一个分支 C 使得 (0, 1) \(\notin \overline{\mathbf{C}}\){[}与习题 4 a) 及 15 b) 相比较{]}。*

\(^{*}\) 17) 设 \(I_{0}\) 为 R 中的区间 \([-1,1]\)，设 \(P_{0}\) 为把 \(I_{0}\) 的两点 1/2 与 1 等同起来所得的商空间；设 \(B_{0}\) 为把 \(I_{0}\) 的点 1/2 与 1 等同且把 -1/2 与 -1 等同所得的商空间。设 I 为 R 中的开区间 {]}--- I, I{[}，并设 P（相应地，B）为 I 在 \(P_{0}\)（相应地，\(B_{0}\)）中的典范像。

\begin{enumerate}
\def\labelenumi{\alph{enumi})}
\item
  证明：三个空间 I、P、B 中任何两个都不同胚（为证 I 与 P 不同胚，注意 P 中有些点，其余集是连通的；对其他空间对亦然）。
\item
  设 X 为可数无限个同胚于 I 的空间与可数无限个同胚于 B 的空间的和。设 Y 为 X 与 P 的和。证明：X 与 Y 不同胚，但 (i) 存在 X 到 Y 上的连续双射及 Y 到 X 上的连续双射；(ii) X 同胚于 Y 的一个开子空间，且 Y 同胚于 X 的一个开子空间。*
\end{enumerate}

*18) 设 X 为实平面 \(\mathbf{R}^2\)，设 R 为 X 上的等价关系，其等价类是诸直线 \(y = \beta\)（对 \(\beta > 0\)）以及诸半直线 \(x = \alpha\)，\(y \leqslant 0\)（对一切 \(\alpha \in \mathbb{R}\)）。证明：若把 X/R 视为 \(\mathfrak{P}_0(X)\) 的一个子集，则 X/R 上由 \(\mathcal{C}_{\Omega}\)（ \(\S 2\) ，习题 7）诱导的拓扑是离散的，而 X/R 上由 \(\mathcal{C}_{\Phi}\) 诱导的拓扑不是离散的，但使 X/R 成为非连通空间。最后证明：X/R 在商拓扑下是连通的。*

*¶ 19) 设 X 为局部紧空间。可以证明，X 可等同于在 § 9 习题 25 中构造的泛紧空间 \(\tilde{X}\) 的一个稠密开子空间（参见第九章，§ 1，习题 8）。再设 X 是连通且局部连通的。

\begin{enumerate}
\def\labelenumi{\alph{enumi})}
\item
  证明：对 X 的每个紧子集 K，X-K 的非相对紧的分支所成之集是有限的{[}参见习题 14 b){]}。由此推出：\(\tilde{X} - X\) 的一个点至多属于 X-K 的分支中一个的闭包{[}若 \(A_i\)（ \(i \leqslant i \leqslant n\)）是 X-K 的非紧分支，而 \(C_i\) 是 \(\tilde{X} - X\) 中位于 \(A_i\) 的闭包中的点所成的紧集，试构造一个以 X 与诸集 \(C_i\) 的和为底集的紧空间，使得 X 等同于该空间的一个稠密开子空间{]}。
\item
  设 \(R'\{x, y\}\) 为 \(\tilde{X}-X\) 上的下述等价关系：“对 X 的每个紧子集 K，x 与 y 位于 X-K 的同一个非紧分支的闭包中”。设 R 为 \(\tilde{X}\) 上的等价关系，其等价类是诸 mod \(R'\) 等价类以及由 X 的单个点组成的子集。证明：R 是豪斯多夫的，且 X 可等同于紧空间 \(X^{b}=\tilde{X}/R\) 的一个稠密开子空间；并证明：在紧空间 \(X^{b}-X\) 中，每个点都是它的既开又闭的邻域之交。\(X^{b}-X\) 的点称为空间 X 的端。
\item
  设 Y 为紧空间，使得 X 可等同于 Y 的一个稠密开子空间，且在紧空间 Y-X 中每个点都是它的既开又闭的邻域之交 (*)。证明：存在 \(X^{b}\) 到 Y 中的一个连续映射，它在 X 上诱导恒同映射。{[}若 Y-X 是两个不相交的非空闭集 A 与 B 的并，证明存在 Y 中两个不相交的开子集 U、V，使得 \(A \subset U\)、\(B \subset V\)，且使得 \(Y - (U \cup V)\) 是紧的{]}。
\item
  证明：若 X 是 σ-紧的，则 \(X^{b}-X\) 有一个可数基。
\item
  实直线有两个端，而 \(R^{n} (n \geqslant 2)\) 有一个。给出 \(R^{2}\) 的一个连通、局部连通、局部紧的子空间的例子，使其端所成之集具有连续统的势。*
\end{enumerate}

\begin{enumerate}
\def\labelenumi{\arabic{enumi})}
\setcounter{enumi}{19}
\tightlist
\item
  设 X 为豪斯多夫空间，其中每个点都有在 X 中既开又闭的邻域组成的基本系。设 Y 为 X 的子空间，A 为 Y 的一个在 Y 中既开又闭的紧子集；证明：存在 X 的一个在 X 中既开又闭的子集 B，使得 \(B \cap Y = A\)。
\end{enumerate}

¶ 21) a) 证明：拓扑空间 X 的下列条件等价：(i) 对 X 的每个开集 U，\(\overline{U}\) 是开的；(ii) 对 X 的每个闭集 F，\(\dot{F}\) 是闭的；(iii) 对 X 中每对不相交的开集 U、V，有 \(\overline{U} \cap \overline{V} = \varnothing\)。满足这些条件的豪斯多夫空间称为极不连通的；它这时是全不连通的。*有理直线是全不连通的但不是极不连通的。*

\begin{enumerate}
\def\labelenumi{\alph{enumi})}
\setcounter{enumi}{1}
\item
  豪斯多夫空间是极不连通的，当且仅当相伴随的半正则空间（§ 8，习题 20）是极不连通的。
\item
  设 X 为豪斯多夫空间，A 为 X 的稠密子集。证明：若 A 是极不连通的，且 X 的拓扑是 A-极大的（§ 3，习题 11），则 X 是极不连通的。
\item
  由 b) 与 c) 推出：若 X 是任一无限集，则 X 的（紧）超滤子空间是极不连通的（参见§ 9，习题 26 与 24）(**)。
\item
  设 X 为没有孤立点的豪斯多夫空间。证明：X 的拓扑是拟极大的（§ 2，习题 6），当且仅当 X 是次极大的（§ 8，习题 22）且极不连通的。（为证条件充分，利用 § 8 习题 22 e)，并证明 X 的每个没有孤立点的闭子集都是开的）。
\item
  证明：每个超正则空间（§ 8，习题 25）都是极不连通的，但没有孤立点的极不连通紧空间不是超正则的（参见§ 9，习题 27）。
\item
  证明：极不连通的半正则空间（§ 8，习题 20）是正则的（参见第二章，§ 4，习题 12）。
\item
  证明：在极不连通空间中，不存在具有无限多个不同项的收敛序列{[}用反证法，递归地构造两个两两不相交的开集序列 \((\mathbf{U}_{n})\)、\((\mathbf{V}_{n})\)，使得若 \(U = \bigcup_{n} U_{n}\) 且 \(V = \bigcup_{n} V_{n}\)，则 \(U \cap V = \varnothing\) 但 \(\overline{U} \cap \overline{V} \neq \varnothing\){]}。
\item
  证明：在极不连通空间 X 中，非孤立点 a 不能有同时既开又闭的、（关于关系 \(\supset\)）良序化的邻域基本系。（如 h) 中那样用反证法。）
\end{enumerate}

¶ 22) a) 证明：在极不连通空间中，每个开子空间与每个稠密子空间都是极不连通的。

\begin{enumerate}
\def\labelenumi{\alph{enumi})}
\setcounter{enumi}{1}
\item
  设 X 为豪斯多夫空间，其中每个点都有一个作为 X 的极不连通子空间的邻域。证明：X 是极不连通的。
\item
  给出一个豪斯多夫空间 X 的例子，它不是极不连通的，但有一个极不连通的稠密开子空间（把两个极不连通空间粘合起来）。
\item
  设 X 为可数离散空间，\(\tilde{X}\) 为 X 上的超滤子所成的（紧）空间（§ 9，习题 26）。设 \((\mathbf{A}_{n})_{n\in\mathbb{N}}\) 是把 X 分成无限子集的一个可数无限划分；在闭子空间 \(Y=\tilde{X}-X\)（\(\tilde{X}\) 的）中，诸集 \(B_{n}=\tilde{X}_{n}\cap Y\) 是开的、闭的且两两不相交的。设 \(B=\bigcup_{n}B_{n}\)；证明：B 在 Y 中的闭包 B 在 Y 中不是开的，从而 Y 不是极不连通的 (*)。（用反证法：利用习题 20，设 C 是 \(\tilde{X}\) 的一个既开又闭的子集，使得 \(C\cap Y=B\)；考察点 \(x_{n}\)（诸集 \(C\cap A_{n}\) 的每一个中的），设 \(\bar{J}\) 为诸点 \(x_{n}\) 所成之集 J 的闭包，并证明 \(\bar{J}\) 不与 B 相交）。
\end{enumerate}

\begin{enumerate}
\def\labelenumi{\arabic{enumi})}
\setcounter{enumi}{22}
\item
  给出双射映射 \(f\)（豪斯多夫空间 \(\mathbf{X}\) 到豪斯多夫空间 \(\mathbf{Y}\) 上的）的一个例子，使得在 \(f\) 下，\(\mathbf{X}\) 的每个紧子集的像是 \(\mathbf{Y}\) 的一个紧子集，且在 \(f\) 下，\(\mathbf{X}\) 的每个连通子集的像是 \(\mathbf{Y}\) 的一个连通子集，但 \(f\) 不是连续的（参见§ 9，习题 4）。
\item
  设 X 为拓扑空间。
\end{enumerate}

\begin{enumerate}
\def\labelenumi{\alph{enumi})}
\tightlist
\item
  设集 \(\mathfrak{P}_{\Omega}(\mathbf{X})\) 赋有诸拓扑 \(\mathfrak{C}_{\Omega}, \mathfrak{C}_{\Phi}\)（ \(\S 2\) ，习题 7）或 \(\mathfrak{C}_{\Theta}\)（ \(\S 8\) ，习题 12）之一。证明：若子集 \(\mathfrak{B}\)（\(\mathfrak{P}_{\Omega}(\mathbf{X})\) 的）是连
\end{enumerate}

通的，且每个 \(\mathbf{M} \in \mathfrak{B}\) 都是连通的，则 \(\bigcup_{\mathbf{M} \in \mathfrak{M}} \mathbf{M}\) 是连通的。

\begin{enumerate}
\def\labelenumi{\alph{enumi})}
\setcounter{enumi}{1}
\tightlist
\item
  设 \(\mathfrak{S}\) 为 \(\mathfrak{P}_{0}(\mathbf{X})\) 的一个子集，它包含 \(\mathbf{X}\) 的所有有限非空子集所成之集。证明：若 \(\mathbf{X}\) 是连通的，则 \(\mathfrak{S}\) 也是连通的{[}利用 § 2 习题 7 b) 与 § 4 命题 8，并考察诸映射 \((x_{1},\ldots ,x_{n})\to \{x_{1},\ldots ,x_{n}\} ]\)。
\end{enumerate}

\begin{enumerate}
\def\labelenumi{\arabic{enumi})}
\setcounter{enumi}{24}
\tightlist
\item
  给定两个拓扑空间 X、Y，映射 \(f: \mathbf{X} \to \mathbf{Y}\) 称为局部同胚，如果对每个 \(x \in \mathbf{X}\) 存在 \(x\) 的一个邻域 U，使得 \(f(\mathbf{U})\) 是 \(f(x)\) 在 Y 中的一个邻域，且映射 \(\mathbf{U} \to f(\mathbf{U})\)（在 U 上与 \(f\) 相合）是一个同胚。每个局部同胚都是连续的开映射。
\end{enumerate}

\begin{enumerate}
\def\labelenumi{\alph{enumi})}
\item
  设 X 是豪斯多夫的，Y 是局部连通的。设 \(f: \mathbf{X} \to \mathbf{Y}\) 为局部同胚，具有这样的性质：存在整数 \(n > 0\)，使得 \(\overline{f}^1(y)\) 恰有 \(n\) 个点（对每个 \(y \in \mathbf{Y}\)）。证明：每个 \(y \in \mathbf{Y}\) 都有一个开邻域 V，使得 \(\overline{f}^1(V)\) 有 \(n\) 个分支 \(\mathbf{U}_i\)（ \(1 \leqslant i \leqslant n\)），且映射 \(\mathbf{U}_i \to \mathbf{V}\)（它与 \(f\) 在 \(\mathbf{U}_i\) 上相合）是同胚。这时 \(f\) 是逆紧映射。
\item
  设 X 是豪斯多夫的，且 Y 是连通且局部连通的。设 \(f: \mathbf{X} \to \mathbf{Y}\) 为逆紧的局部同胚。证明：\(f\) 满足 \(a\) ) 的假设。{[}若对每个整数 \(n > 0\)，\(\mathbf{Y}_n\) 是这样的 \(y \in \mathbf{Y}\) 所成之集：\(\bar{f}^1(y)\) 至少有 \(n\) 个点，证明 \(\mathbf{Y}_n\) 在 \(\mathbf{Y}\) 中既开又闭{]}。
\end{enumerate}

\(^{*}c)\) 给出满射局部同胚 \(f: X \to Y\) 的一个例子，其中 Y 是紧的、连通的且局部连通的，X 是局部紧的、连通的且局部连通的，\(\overline{f}^{1}(y)\) 对每个 \(y \in Y\) 至多由两个点组成，但 f 不是逆紧的（参见§ 5，习题 4）。
